\documentclass{book}
\usepackage{hyperref}
\usepackage{amsmath}
\usepackage{ulem}
\usepackage{tipa} % Add this line to support \textipa command
\hypersetup{
    colorlinks=true,
    linkcolor=black,
    filecolor=magenta,      
    urlcolor=cyan,
    }

\urlstyle{same}

\title{French Grammar}
\author{Junzhang Luo}

\begin{document}

\tableofcontents

\chapter{La Phrase Simple}
    The structure of a simple sentence in French is as follows:
    \begin{center}
        Subject + conjugated verb + article + noun
    \end{center}
    For example:
    \begin{center}
        Tu regardes un film (You are watching a movie)
    \end{center}

\chapter{Les Articles Indéfinis}
    \begin{itemize}
        \item Masculine singular: un
        \item Feminine singular: une
        \item Masculine / Feminine plural: des
    \end{itemize}

\chapter{Le Genre de Nom}
    \begin{enumerate}
        \item The obvious genders: un garçon / une fille, un oncle / une tante \dots
        \item 90\% rule: words ending in -E are feminine.
    \end{enumerate}

    \section{Masculine Nouns}
        \begin{enumerate}
            \item -ACLE: un obstacle (an obstacle), un miracle (a miracle) \dots
            \item -AGE: un paysage (a landscape), un coloriage (a coloring) \dots
            \begin{itemize}
                \item Exceptions: une cage, une image, une plage
            \end{itemize}
            \item -AL: un cheval (a horse), un canal (a canal), un journal (a newspaper) \dots
            \item -EAU: un morceau (a piece), un chapeau (a hat), un manteau (a coat) \dots
            \begin{itemize}
                \item Exceptions: une peau (a skin), une eau (a water)
            \end{itemize}
            \item -EU: un cheveu (a hair), un jeu (a game), un vœu (a wish) \dots
            \item -ET: un billet (a ticket), un objet (an object), un secret (a secret) \dots
            \item -IER: un pommier (an apple tree), un panier (a basket) \dots
            \item -IN: un cousin (a cousin), un lapin (a rabbit), un matin (a morning) \dots
            \item -ISNE: un voisin (a neighbor), un cousin (a cousin) \dots
            \item -MENT: un appartement (an apartment), un vêtement (a clothing) \dots
            \item -OIR: un miroir (a mirror), un soir (an evening) \dots
            \item -OU: un clou (a nail), un genou (a knee) \dots
            \item -ON: un avion (an airplane) \dots
            \begin{itemize}
                \item Exceptions: une prison (a prison) \dots
            \end{itemize}
        \end{enumerate}
    
    \section{Masculine only Topics}
        \begin{enumerate}
            \item Days of the week: le lundi \dots
            \item Seasons: un printemps, un été, un automne, un hiver
            \item Languages: Le grecque, le français, le chinois \dots
            \begin{itemize}
                \item Languages are not capitalized in French.
            \end{itemize}
            \item Weight \& Metrics: un kilomètre, un gramme, un litre \dots
            \item English words: un week-end, un email, un parking \dots
        \end{enumerate}
    
    \section{Feminine Nouns}
        \begin{enumerate}
            \item -ADE: une balade (a walk), une cascade (a waterfall) \dots
            \item -ALE: une rivale (a rival), une baleine (a whale) \dots
            \item -ANCE: une chance (a chance), une balance (a scale) \dots
            \item -ENCE: une exigence (a requirement), une absence (an absence) \dots
            \item -ETTE: une vedette (a star), une fourchette (a fork) \dots
            \item -IE: une galaxie (a galaxy), une harmonie (a harmony) \dots
            \item -IQUE: une brique (a brick), une musique (a music) \dots
            \item -OIRE: une viotoire (a victory), une histoire (a story) \dots
            \item -SION: une passion (a passion), une discussion (a discussion) \dots
            \item -TÉ / TIÉ: une société (a society), une amitié (a friendship) \dots
            \begin{itemize}
                \item Exception: un été (a summer)
            \end{itemize}
            \item -TION: une parition (a publication), une question (a question) \dots
            \item -URE: une auenture (an adventure), une culture (a culture) \dots
        \end{enumerate}

    \section{Professions \& People}
        \subsection{Only Change Articles in Front}
            Un / une architecte, un / une collègue, un / une enfant, un / une élève 
        \subsection{VOWELS + E}
            This is the most common case, pronunciation does not change.
            \begin{center}
                VOWEL + E
            \end{center}
            Example: un apprenti (an apprentice) / une apprentie (an apprentice)
        \subsection{Consonants + E}
            Pronunciation changes, the different cases are:
            \begin{center}
                \begin{tabular}{|c|c|c|c|}
                    \hline
                    \textbf{Ending} & \textbf{Masculine} & \textbf{Feminine} & \textbf{English}\\
                    \hline
                    -T & un avocat & une avocate & lawyer\\
                    \hline
                    -N & un cousin & une cousine & cousin\\
                    \hline
                    -S & un Anglais & une Anglaise & Englishman \\
                    \hline
                    -ER / -ÈRE & un infirmier & une infirmière & nurse \\
                    \hline
                    -EUR / -EUSE & un coiffeur & une coiffeuse & hairdresser\\
                    \hline
                    -IEN / -IENNE & un pharmacien & une pharmacienne & pharmacist\\
                    \hline
                    -ON / -ONNE & un patron & une patronne & boss\\
                    \hline
                    -TEUR / -TRICE & un acteur & une actrice & actor\\
                    \hline
                    -F / -VE & un veuf & une veuve & widower / widow\\
                    \hline
                    -X / -SE & un époux & une épouse & husband / wife\\
                    \hline
                \end{tabular}
            \end{center}
        
    \section{Countries}
        \begin{itemize}
            \item For country names, use 90\% rule, except \textit{Le Mexique} and \textit{Le Cambodge}.
            \item Nationality example:
                \begin{itemize}
                    \item L'Algérie (f): Algérien (m) / Algérienne (f)
                \end{itemize}
        \end{itemize}
        
    \section{Nouns with Specific Gender}
    These nouns does not care about the person's gender, they are always the same.
        \begin{itemize}
            \item un bébé (a baby)
            \item une star (a star)
            \item une vedette (a star)
            \item une victime (a victim)
            \item \dots
        \end{itemize}
    
    \section{Same Noun Different Meanings}
        \begin{itemize}
            \item Le mort (dead person) / la mort (death)
            \item Le chèvre (goat cheese) / la chèvre (goat)
            \item Le poêle (stove) / la poêle (frying pan)
            \item Le vase (vase) / la vase (mud in water)
            \item \dots
        \end{itemize}
    
    \section{Le Pluriel des Nouns}
        \begin{enumerate}
            \item -E: add -S
            \item -S / -X / -Z: no change
            \item -EAU / -AU / -EU: add -X for most cases
            \begin{itemize}
                \item Exception: 
                \begin{itemize}
                    \item un landau / des landous (a baby carriage)
                    \item un neveu / des neveux (a nephew)
                    \item un bleu / des bleus (a bruise)
                \end{itemize}
            \end{itemize}
            \item -OU: regular takes -s, irregular takes -x, irregular is the majority.
            \begin{itemize}
                \item -S Examples:
                \begin{itemize}
                    \item un bisou, des bisous (a kiss)
                    \item un clou / des clous (a nail)
                \end{itemize}
                \item -X Examples:
                \begin{itemize}
                    \item un bijou / des bijoux (a jewel)
                    \item un caillou / des cailloux (a pebble)
                    \item un chou / des choux (a cabbage)
                    \item un genou / des genoux (a knee)
                    \item un hibou / des hiboux (an owl)
                    \item un joujou / des joujoux (a toy)
                    \item un pou / des poux (a louse)
                \end{itemize}
                * The phrase to memorize is: 
                \begin{center}
                    Viens mon \textcolor{red}{chou}, mon \textcolor{red}{hibou}, mon \textcolor{red}{bijou}, mon \textcolor{red}{joujou}, sur mes genoux, et jette des \textcolor{red}{cailloux} a` ce \textcolor{red}{hibou} plein de \textcolor{red}{poux}. \\
                    (Come my cabbage, my owl, my jewel, my toy, on my knees, and throw pebbles at this owl full of lice.)
                \end{center}
            \end{itemize}
            \item -AL: 
            \begin{enumerate}
                \item The common case is -AUX: un animal / des animaux \dots
                \item The uncommon case is -ALS (about festivity): un bal / des bals, un carnaval / des carnavals \dots
            \end{enumerate}
            \item -AIL:
            \begin{enumerate}
                \item -AILS:
                \begin{center}
                    \begin{tabular}{|c|c|c|}
                        \hline
                        \textbf{Singular} & \textbf{Plural} & \textbf{Meaning}\\
                        \hline
                        un chandail & des chandails & a sweater\\
                        \hline
                        un détail & des détails & a detail\\
                        \hline
                        un épouvantail & des épouvantails & a scarecrow\\
                        \hline
                        \dots & \dots & \dots\\
                        \hline
                    \end{tabular}
                \end{center}
                \item -AUX:
                \begin{center}
                    \begin{tabular}{|c|c|c|}
                        \hline
                        \textbf{Singular} & \textbf{Plural} & \textbf{Meaning}\\
                        \hline
                        un bail & des baux & a lease\\
                        \hline
                        un corail & des coraux & a coral\\
                        \hline
                        un émail & des émaux & an enamel\\
                        \hline
                        \dots & \dots & \dots\\
                        \hline
                    \end{tabular}
                \end{center}
            \end{enumerate}
            \item Ending with consonnat: add -S
            \item -ON / -AN / -IN / -UM: add -S
        \end{enumerate}
    
    \section{Irregular Plurals}
        \begin{center}
            \begin{tabular}{|c|c|c|}
                \hline
                \textbf{Singular} & \textbf{Plural} & \textbf{Meaning}\\
                \hline
                un œil & des yeux & an eye\\
                \hline
                un ciel & des cieux & a sky\\
                \hline
                madame / mademoiselle & mesdames / mesdemoiselles & Mrs. / Miss\\
                \hline
            \end{tabular}
        \end{center}
        Pronounciation changes:
        \begin{itemize}
            \item un œil / des yeux
            \item un boeuf / des boeufs 
        \end{itemize}
    
    \section{Nouns that are Always Plural}
        Les vacances, les mathématiques, les lunettes, les gens \dots
    
    \section{Singular but Plural}
        Le public, La foule (the crowd), La famille, Le groupe \dots

\chapter{Les Pronoun Subjects}

    \section{Definite Articles (Le / La / Les)}
        The following are additional rules on when to use definite articles in French that are different from English:
        \begin{enumerate}
            \item For abstract nouns, i.e. nouns that are not physical objects.
            \item For concrete nouns when talking about a general concept.
            \begin{center}
                Le café contient de la caféine. (Coffee contains caffeine.)
            \end{center}
            \item For language and school subject.
            \item For continent, countries, provinces and regions.
            \item With quantities and prices.
            \begin{center}
                50 kilomètres à l'heure, 2 euros la pièce
            \end{center}
            \item Dates \& days of the week.
                \begin{itemize}
                    \item For days of week, only use for things that happen regularly, add \textit{le} in front.
                    \begin{center}
                        Elle travaille le lundi (She works on Monday) \\
                        Elle travaille lundi (She works next Monday)
                    \end{center}
                \end{itemize}
            \item Body parts, when referring to personal body parts, do not use 'my \dots' in French.
            \begin{center}
                Je me brosse les cheveux (I brush my hair) \\
                Il s'est cassé le bras (He broke his arm)
            \end{center}
            \item For \textit{parle le \dots}, \textit{le} is optional.
            \begin{itemize}
                \item For H aspiré, use \textit{le}.
                \item For H muet, use \textit{l'}.
            \end{itemize}
        \end{enumerate}

    \section{Les Articles Partitifs}
        Use to talk about smoething without specifying the quantity.
        \begin{enumerate}
            \item Masculine Singular: 
            \begin{center}
                du (de + le)
            \end{center}
            use for masculine noun starting with a consonant.
            \item Feminine Singular: de la
            \item Masculine \& Feminine Singular starting with a vowel: de l'
            \item Pluriel: des
            \item *Always use \textit{de} in negative sentences.
        \end{enumerate}

    \section{Conjugation of Être}
        \begin{center}
            \begin{tabular}{|c|c|}
                \hline
                \textbf{Subject} & \textbf{Conjugation}\\
                \hline
                Je & suis\\
                \hline
                Tu & es\\
                \hline
                Il / Elle / On & est\\  
                \hline
                Nous & sommes\\
                \hline
                Vous & êtes\\
                \hline
                Ils / Elles & sont\\
                \hline
            \end{tabular}
        \end{center}

    \section{C'est / Ce sont}
        \begin{enumerate}
            \item For persons, things, nationalities and jobs:
                \begin{center}
                    C'est / Ce sont + article + masculine / feminine noun
                \end{center}
            \item For opinions:
                \begin{center}
                    C'est + adjective / adverb 
                \end{center}
                always sigular. For speaking French, usually use \textit{C'est} for plural.
            \item \textit{Il est}:
                \begin{enumerate}
                    \item Description: 
                        \begin{center}
                            Il est  + adjective / adverb
                        \end{center}
                    \item Time:
                        \begin{center}
                            Il est + time
                        \end{center}
                    \item Jobs:
                        \begin{center}
                            Il est + masculine noun
                        \end{center}
                    \item Nationalities:
                        \begin{center}
                            Il est \dots
                        \end{center}
                        Ex. 
                        \begin{center}
                            Il est suisse (He is Swiss) \\
                            C'est un Suisse (He is a Swiss)
                        \end{center}
                        \textbf{Note.} that the nationality is a noun in the second sentence, so it is capitalized.
                \end{enumerate}
        \end{enumerate}

    \section{Les Adjectifs Démonstratifs}
        \begin{enumerate}
            \item Masculine singular: Ce / Cet
            \item Feminine singular: Cette
            \item Plural: Ces
        \end{enumerate}

    \section{Les Adjectifs Possessifs}
        \begin{center}
            \begin{tabular}{|c|c|c|c|}
                \hline
                \textbf{Subject} & \textbf{Masculine Singular} & \textbf{Feminine Singular} & \textbf{Plural}\\
                \hline
                Je & mon & ma & mes\\
                \hline
                Tu & ton & ta & tes\\
                \hline
                Il / Elle / On & son & sa & ses\\
                \hline
                Nous & notre & notre & nos\\
                \hline
                Vous & votre & votre & vos\\
                \hline
                Ils / Elles & leur & leur & leurs\\
                \hline
            \end{tabular}
        \end{center}

\chapter{Compter}
    \begin{center}
        \begin{tabular}{|c|c|c|c|}
            \hline
            0: Zéro & 10: Dix & 20: Vingt & 21: Vingt et un\\
            \hline
            1: Un & 11: Onze & 30: Trente & 22: Vingt-deux \\
            \hline
            2: Deux & 12: Douze & 40: Quarante & 23: Vingt-trois \\
            \hline
            3: Trois & 13: Treize & 50: Cinquante & 24: Vingt-quatre \\
            \hline
            4: Quatre & 14: Quatorze & 60: Soixante & 25: Vingt-cinq \\
            \hline
            5: Cinq & 15: Quinze & 70: Soixante-dix & 26: Vingt-six \\
            \hline
            6: Six & 16: Seize & 80: Quatre-vingts & 27: Vingt-sept \\
            \hline
            7: Sept & 17: Dix-sept & 90: Quatre-vingt-dix & 28: Vingt-huit \\
            \hline
            8: Huit & 18: Dix-huit & 100: Cent & 29: Vingt-neuf \\
            \hline
            9: Neuf & 19: Dix-neuf & & \\
            \hline
        \end{tabular}
    \end{center}
    \begin{enumerate}
        \item For 20 to 69: let $x \in \left\{20, 30, \dots, 60\right\}$, $y \in \left\{ 2, \dots, 9\right\}$, then the format is:
                $$ [x-y] $$
                $x1$ is $x$ \textit{et un}
        \item For 70: (60 + 10), $7x$ is (60 + $1x$), $x \in \left\{1, \dots, 9\right\}$.
        \item For 80: (4 * 20), $8x$ is (4 * $20 + x$), $x \in \left\{0, \dots, 19\right\}$.
        \item For 90: (4 * 20 + 10), $9x$ is (4 * $20 + 10 + x$), $x \in \left\{1, \dots, 9\right\}$.
    \end{enumerate}

    \section{Compter de 100 à 1000}
        \begin{itemize}
            \item 1000: Mille
            \item Except for 100, when $x \equiv 0 \mod 100$, then $x$ \textit{cents} 
        \end{itemize}
    
    \section{Compter Après 1000}
        \begin{itemize}
            \item \textit{Mille} never has -s in the end.
            \item \textit{Million, Billion, Trillion} takes -s in the end except for 1.
        \end{itemize}
    
    \section{Ordinal Numbers}
        \begin{center}
            number + -ième \\
            \begin{tabular}{|c|c|c|}
                \hline
                un / une & premier / première & \\
                \hline
                deux & deuxième & \\
                \hline
                trois & troisième & \\
                \hline
                quatr\textcolor{red}{e} & quatrième & drop the -e \\
                \hline
                cinq & cinq\textcolor{red}{u}ième & add -u \\
                \hline
                six & sixième & \\
                \hline
                sept & septième & \\
                \hline
                huit & huitième & \\
                \hline
                neuf & neuvième & -f $\to$ -v \\
                \hline
                dix & dixième & \\
                \hline
            \end{tabular}
        \end{center}
        In general, drop the -e. Ex. \textit{trente} $\to$ \textit{trentième}.

\chapter{Negation}
    The negation in French is as follows:
    \begin{center}
        ne + verb + pas
    \end{center}
    \textit{ne \dots pas} can also be replaced with the following:
    \begin{enumerate}
        \item ne \dots plus (no more)
        \item ne \dots jamais (never)
        \item ne \dots que (only)
        \item ne \dots rien (nothing)
        \item ne \dots personne (nobody)
        \item ne \dots pas encore (not yet)
        \item ne \dots ni \dots ni \dots (neither \dots nor \dots)
    \end{enumerate}
    Rules for negation:
    \begin{itemize}
        \item If any of \textit{un, une, des, du, de la, de l', des} follows the negation, use \textit{de / d'} instead.
        \item If the negation is followed by \textit{le, la, l', les, c'est / être}, then the article does not change.
    \end{itemize}

    \section{Negation in Spoken French}
        \textit{Ne} can be dropped in informal spoken French.

    \section{Negation with an Infinitive Verb}
        Infinitive verb is the verb that is not conjugated. The negation is as follows:
        \begin{center}
            ne pas + infinitive verb
        \end{center}
        For example:
        \begin{center}
            Il me demande de ne pas venir maintenant (He asks me not to come now) \\
            Merci de ne pas fumer (Thank you for not smoking) \\
            Ne pas déranger (Do not disturb) \\
            Ne pas marcher sur la pelouse (Do not walk on the grass)
        \end{center}

    \section{Passé Composé and Negation}
        \begin{center}
            ne / n' + avoir / être + pas + past participle
        \end{center}
        \begin{enumerate}
            \item Normal case: Je n'ai pas vu le chien (I did not see the dog)
            \item Special case: Vous n'avez reccontré personne (You did not meet anyone)
        \end{enumerate}

    \section{Passé Composé Negation of Reflexive Verbs}
        \begin{center}
            ne / n' + reflexive pronoun + être + pas + past participle
        \end{center}
    
\chapter{Il y a}
    Always masculine singular, use \textit{il y a} for both singular and plural.
    \begin{center}
        Il y a = avoir (to have)
    \end{center}
    \textit{a} can be conjugated in other tenses: \textit{avait} (was), \textit{aura} (will have) \dots
    \begin{enumerate}
        \item Followed by masculine / feminine noun: Il y a \dots
        \item In negative sentence: Il n'y a pas de / d' \dots
        \item Il y a = ago: \dots il y a [time]
    \end{enumerate}

\chapter{Préposition}

    \section{La Préposition à} 
        \begin{enumerate}
            \item Translate to \textit{to, at, in, for from, per} in English.
            \item Changes bby the article followed after:
                \begin{enumerate}
                    \item à + le = au + masculine singular noun
                    \item à + la = à la + feminine singular noun
                    \item à + l' = à l' + noun starting with a vowel / h muet
                    \item à + les = aux + plural noun
                \end{enumerate}
        \end{enumerate}
    
        \subsection{Quand l'utiliser}
            \begin{enumerate}
                \item Location: à + city
                    \begin{center}
                        Je vais à Paris (I am going to Paris)
                    \end{center}
                \item Distance, time, age, rate: à + \#
                    \begin{center}
                        Il habite à 5 kilomètres (He lives 5 kilometers away)
                    \end{center}
                \item Manner \& characteristic
                    \begin{center}
                        Il parle à voix basse (He speaks in a low voice)
                    \end{center}
                \item Procession
                    \begin{center}
                        un ami à moi (a friend of mine)
                    \end{center}
                \item Food and how its made
                    \begin{center}
                        une tarte aux pommes (an apple pie)
                    \end{center}
                \item Specific transportation: à pied, à cheval, à vélo \dots
                \item To express pain:
                    \begin{center}
                        avoir mal à la tête (to have a headache)
                    \end{center}
                \item Name of feastivals
                \item To show purpose of smoething
                    \begin{center}
                        un couteau à beurre (a butter knife)
                    \end{center}
                    \textbf{Note.} à stays as à.
                \item Say goodbye
                    \begin{center}
                        à demain (see you tomorrow)
                    \end{center}
                \item Few expression with \textit{Être}
                    \begin{itemize}
                        \item Être a` qqn (to belong to someone)
                        \item Être au courant (to be aware)
                        \item Être à l'heure (to be on time)
                    \end{itemize}
                \item C'est and à
                    \begin{itemize}
                        \item C'est bon à savoir (It's good to know)
                        \item C'est difficile à croire (It's hard to believe)
                    \end{itemize}
                \item Verbs + à:
                    \begin{enumerate}
                        \item Verb + à + noun:
                            \begin{center}
                                \begin{tabular}{|c|c|}
                                    \hline
                                    \textbf{Verb} & \textbf{Meaning}\\
                                    \hline
                                    Acheter qqch à qqn & to buy something for someone\\
                                    Demander qqch à qqn & to ask someone for something\\
                                    Envoyer qqch à qqn & to send something to someone\\
                                    Faire mal à qqn & to hurt someone\\
                                    Jour à qqch & to play something\\
                                    Plaire à qqn & to please someone\\
                                    Répondre à qqn & to answer someone\\
                                    Téléphoner à qqn & to call someone\\
                                    \dots & \dots\\
                                    \hline
                                \end{tabular}
                            \end{center}
                        
                        \item Verbs + à + infinitive verb:
                            \begin{center}
                                \begin{tabular}{|c|c|}
                                    \hline
                                    \textbf{Verb} & \textbf{Meaning}\\
                                    \hline
                                    Aider à faire qqch & to help to do something \\
                                    Chercher à faire qqch & to attempt to do something \\
                                    Commencer à faire qqch & to begin to do something \\
                                    Hésiter à faire qqch & to hesitate to do something \\
                                    Se mettre à faire qqch & to start doing something \\
                                    Réfléchir à faire qqch & to think of doing something \\
                                    Servir à faire qqch & to be used to do something \\
                                    Tenir à faire qqch  & to insist on doing something \\
                                    \dots & \dots \\
                                    \hline
                                \end{tabular}
                            \end{center}
                    \end{enumerate}
                
                \item Être + adjective + à:
                    \begin{itemize}
                        \item Use to describe the state or feeling of someone. \\
                                Ex. Être prêt(e) à (to be ready to)
                    \end{itemize}
                
                \item Le (la) seul(e) + à (only one(s))
                \item Nombreux(euses) + à (many is / are)
                \item Ordinal number + à \& le dernier à (the last one)
                    
            \end{enumerate}

    \section{Le Préposition de}
        \begin{enumerate}
            \item Translate to \textit{from, of, about} in English.
            \item Changes bby the article followed after:
                \begin{enumerate}
                    \item de + le = du + masculine singular noun
                    \item de + la = de la + feminine singular noun
                    \item de + l' = de l' + noun starting with a vowel / h muet
                    \item de + les = des + plural noun
                \end{enumerate}
        \end{enumerate}
    
        \subsection{Quand l'utiliser}
            \begin{enumerate}
                \item Starting point \& origin
                    \begin{center}
                        Il vient de Paris (He comes from Paris)
                    \end{center}
                \item Possession
                    \begin{center}
                        le livre de Marie (Marie's book)
                    \end{center}
                \item Indicate what something is made of
                    \begin{center}
                        une table de bois (a wooden table)
                    \end{center}
                \item From \dots to with à
                    \begin{center}
                        de Paris à Londres (from Paris to London)
                    \end{center}
                \item Superlatives (most / best etc.)
                    \begin{center}
                        le plus grand de tous (the tallest of all)
                    \end{center}
                \item Quantities (stay as 'de')
                    \begin{center}
                        un kilo de pommes (a kilo of apples)
                    \end{center}
                \item Expressions with \textit{Être} and \textit{de}
                    \begin{center}
                        Être de \dots (to be) \\
                        Ex. Être d'accord (to agree)
                    \end{center}
                \item Verbs + de:
                    \begin{enumerate}
                        \item Verb + de + noun:
                            \begin{center}
                                \begin{tabular}{|c|c|}
                                    \hline
                                    \textbf{Verb} & \textbf{Meaning}\\
                                    \hline
                                    S'approcher de qqn / qqch & to approach someone / something\\
                                    Arriver de + endroit & to arrive from + place\\
                                    Avoir besion de qqn / qqch & to need someone / something\\
                                    Parler de qqn / qqch & to talk about someone / something\\
                                    Penser de qqn / qqch & to think of someone / something\\
                                    Réver de qqn / qqch & to dream of someone / something\\
                                    Se souvenir de qqn / qqch & to remember someone / something\\
                                    \dots & \dots\\
                                    \hline
                                \end{tabular}
                            \end{center}
                        
                        \item Verb + de + infinitive verb:
                            \begin{center}
                                \begin{tabular}{|c|c|}
                                    \hline
                                    \textbf{Verb} & \textbf{Meaning}\\
                                    \hline
                                    Accepter de faire & to accept to do something \\
                                    Avoir raison de faire & to be right to do something \\
                                    Conseiller de faire & to advise to do something \\
                                    Offrir de faire & to offer to do something \\
                                    Refuser de faire & to refuse to do something \\
                                    Re2ver de faire & to dream of doing something \\
                                    Se souvenir de faire & to remember doing something \\
                                    Venir de faire & to have just done something \\
                                    \dots & \dots \\
                                    \hline
                                \end{tabular}
                            \end{center}
                    \end{enumerate}
                
                \item Être + en train + de + infinitive verb (to be in the process of doing)
                    \begin{center}
                        Je suis en train d'écrire une lettre (I am writing a letter)
                    \end{center}
                
                \item Être + sur le point + de + infinitive verb (to be about to do)
                    \begin{center}
                        Il est sur le point de partir (He is about to leave)
                    \end{center}
                
                \item Être + adj + de + infinitive verb / noun (talk about feelings) \\
                        Ex. \textit{Être anxieux(e) de} (to be anxious to)
                \item Il est + adj + de + infinitive verb:
                    \begin{itemize}
                        \item The adjective always stay as masculine singular.
                        \item Use to describe the state or feeling of someone.
                        \item Ex. \textit{Il est difficile de} (It is difficult to)
                    \end{itemize}
                \item C'est + adj + de + infinitive verb:
                    \begin{itemize}
                        \item The adjective never changes.
                        \item Ex. \textit{C'est bien de \dots} (It's good to \dots)
                    \end{itemize}  
            \end{enumerate}

    \section{Préposition EN}
        \begin{enumerate}
            \item Translate to \textit{in, on, to} in English.
        \end{enumerate}

        \subsection{Quand l'utiliser}
            \begin{enumerate}
                \item Location \& destination (always use \textit{en} for feminine countries, for masculine countries, use \textit{au})
                    \begin{center}
                        Il est en France (He is in France)
                    \end{center}
                \item Years, month, seasons (with the exception of \textit{Au printemps})
                    \begin{center}
                        en 2020 (in 2020) \\
                        en janvier (in January) \\
                        en été (in summer)
                    \end{center}
                \item The time that something will take
                    \begin{center}
                        Le dîner sera prêt en 10 minutes (Dinner will be ready in 10 minutes)
                    \end{center}
                \item Transportation
                    \begin{center}
                        en avion (by plane) \\
                        en voiture (by car) \\
                        en train (by train) \\
                        \dots
                    \end{center}
                \item Language as a way to say \underline{in}
                    \begin{center}
                        Le formulaire est en anglais (The form is in English)
                    \end{center}
                \item Material
                    \begin{center}
                        Un collier en argent (A silver necklace)
                    \end{center}
                \item Common expressions with \textit{être} and \textit{en}:
                    \begin{center}
                        \begin{tabular}{|c|c|}
                            \hline
                            \textbf{Expression} & \textbf{Meaning}\\
                            \hline
                            Être en avance & to be early\\
                            Être en vacances & to be on vacation\\
                            Être en retard & to be late \\
                            Ëtre en forme & to be in shape \\
                            Être en bonne santé & to be healthy \\
                            Être en panne & to be out of order \\
                            \dots & \dots \\
                            \hline
                        \end{tabular}
                    \end{center}
                    \textbf{Note.} In French, \textit{vacances} is always plural.
            \end{enumerate}

    \section{ALLER + Prépositions}
        Conjugatin of \textit{Aller}:
        \begin{center}
            \begin{tabular}{|c|c|}
                \hline
                Je vais & Nous allons \\
                \hline
                Tu vas & Vous allez \\
                \hline
                Il / Elle / On va & Ils / Elles vont \\
                \hline
            \end{tabular}
        \end{center}
        
        \begin{enumerate}
            \item Aller + à la: followed by feminine noun
                \begin{center}
                    Aller à la plage (to go to the beach)
                \end{center}
            \item Aller + au: followed by masculine noun
                \begin{center}
                    Aller au cinéma (to go to the cinema)
                \end{center}
            \item Aller + aux: followed by plural noun
                \begin{center}
                    Aller aux États-Unis (to go to the United States)
                \end{center}
            \item Aller + en: followed by geography terms (country, province etc.)
                \begin{center}
                    Aller en France (to go to France)
                \end{center}
            \item Aller + au: followed by masculine country
                \begin{center}
                    Aller au Canada (to go to Canada)
                \end{center}
            \item Aller + aux: followed by plural country
                \begin{center}
                    Aller aux Pays-Bas (to go to the Netherlands)
                \end{center}
            \item Aller + à: followed by city
                \begin{center}
                    Aller à Paris (to go to Paris)
                \end{center}
            \item Aller + chez: followed by person, business or organization
                \begin{center}
                    Aller chez le médecin (to go to the doctor)
                \end{center}
        \end{enumerate}

    \section{Les Préposition de Lieu}
        \begin{enumerate}
            \item Contre (against): usually followed by article
                \begin{center}
                    Il est assis contre le mur (He is sitting against the wall)
                \end{center}
            \item Dans (in, inside): usually followed by article
                \begin{center}
                    Il est dans la maison (He is in the house)
                \end{center}
            \item Derrière (behind): usually followed by article
                \begin{center}
                    Il est derrière la maison (He is behind the house)
                \end{center}
            \item Devant (in front of): usually followed by article
                \begin{center}
                    Il est devant la maison (He is in front of the house)
                \end{center}
            \item Entre (between): usually followed by article
                \begin{center}
                    Il est entre les deux maisons (He is between the two houses)
                \end{center}
            \item Sous (under): usually followed by article
                \begin{center}
                    Il est sous la table (He is under the table)
                \end{center}
            \item Sur (on): usually followed by article
                \begin{center}
                    Il est sur la table (He is on the table)
                \end{center}
            \item À côté de (next to): usually followed by article
                \begin{center}
                    Il est à côté de la maison (He is next to the house)
                \end{center}
            \item À droite (gauche) de (to the right (left) of): usually followed by article
                \begin{center}
                    Il est à droite de la maison (He is to the right of the house)
                \end{center}
            \item Au bord de (by / on the edge of): usually followed by article
                \begin{center}
                    Il est au bord de la piscine (He is by the pool)
                \end{center}
            \item Au-dessons de (behind): usually followed by article
                \begin{center}
                    Il est au-dessons de la table (He is behind the table)
                \end{center}
            \item Au-dessus de (above): usually followed by article
                \begin{center}
                    Il est au-dessus de la table (He is above the table)
                \end{center}
            \item Au milieu de (in the middle of): usually followed by article
                \begin{center}
                    Il est au milieu de la table (He is in the middle of the table)
                \end{center}
            \item En face de (facing / in front of): usually followed by article
                \begin{center}
                    Il est en face de la maison (He is facing the house)
                \end{center}
            \item Loin de (far from): usually followed by article
                \begin{center}
                    Il est loin de la maison (He is far from the house)
                \end{center}
            \item Près de (near): usually followed by article
                \begin{center}
                    Il est près de la maison (He is near the house)
                \end{center}
        \end{enumerate}
    
    \section{Les Préposition de Temps}
        \begin{enumerate}
            \item À (indivate exact time):
                \begin{center}
                    Je vais au cinéma à 20 heures (I am going to the cinema at 8 o'clock)
                \end{center}
            \item En (amount of time takes to do something):
                \begin{center}
                    Je vais au cinéma en 10 minutes (I am going to the cinema in 10 minutes)
                \end{center}
            \item Dans (time before something happen):
                \begin{center}
                    Je vais au cinéma dans 10 minutes (I am going to the cinema in 10 minutes)
                \end{center}
            \item Après (after):
                \begin{center}
                    Je vais au cinéma après le dîner (I am going to the cinema after dinner)
                \end{center}
            \item Avant (before):
                \begin{center}
                    Je vais au cinéma avant le dîner (I am going to the cinema before dinner)
                \end{center}
            \item Depuis (since / for): ongoing action used in present tense
                \begin{center}
                    Je suis ici depuis 10 minutes (I have been here for 10 minutes)
                \end{center}
            \item Jusqu'à / au / aux (until, before, to, up to): About duration, Distance
                \begin{center}
                    Je vais au cinéma jusqu'à 20 heures (I am going to the cinema until 8 o'clock)
                \end{center}
            \item Jusqu'à où (distance in questions)
                \begin{center}
                    Jusqu'à où est la plage? (How far is the beach?)
                \end{center}
            \item Jusqu'à en (about time frame from a moment to another, include month and year / about destination):
                \begin{center}
                    Jusqu'en 2020 (Until 2020) \\
                    Jusqu'en France (Until France)
                \end{center}
            \item Durant / Pendant (during): Entire duration of action, specific timeframes, general situations / vague future actions
                \begin{center}
                    Je vais au cinéma pendant le dîner (I am going to the cinema during dinner)
                \end{center}
            \item Pour (for):
                \begin{center}
                    Je vais au cinéma pour 10 minutes (I am going to the cinema for 10 minutes)
                \end{center}
        \end{enumerate}
    
    \section{Les Autres Prépositions}
        \subsection{Par}
            Translate to \textit{through, by, per, for, around} in English. \\
            Used in the following cases:
            \begin{enumerate}
                \item Talk about a reason (by \dots)
                    \begin{center}
                        Il est malade par le froid (He is sick because of the cold)
                    \end{center}
                \item Talk about measurement (per \dots)
                    \begin{center}
                        50 kilomètres par heure (50 kilometers per hour)
                    \end{center}
                \item Directions (by \dots, around \dots)
                    \begin{center}
                        Il est passé par la rue (He went by the street)
                    \end{center}
                \item Verbs + par:
                    \begin{center}
                        \begin{tabular}{|c|c|}
                            \hline
                            \textbf{Verb} & \textbf{Meaning}\\
                            \hline
                            entrer par & to enter by \\
                            finir par & to end up by \\
                            habiter par ici & to live around here \\
                            venir par & to come by \\
                            \dots & \dots \\
                            \hline
                        \end{tabular}
                    \end{center}
            \end{enumerate}
        
        \subsection{Pour}
            Translate to \textit{for, per, in favor of, on behalf of} in English. \\
            Used in the following cases:
            \begin{enumerate}
                \item Destination, people, place, abstract goal
                    \begin{center}
                        Il est parti pour Paris (He left for Paris)
                    \end{center}
                \item Verbs + pour:
                    \begin{center}
                        Ex. Ca compte pour du beurre (It counts for nothing) \\
                        \begin{tabular}{|c|c|}
                            \hline
                            \textbf{Verb} & \textbf{Meaning}\\
                            \hline
                            compter pour & to be worth \\
                            être pour & to be in favor of \\
                            partir pour & to leave for \\
                            votre pour & to vote for \\
                            \dots & \dots \\ 
                            \hline
                        \end{tabular}
                    \end{center}
            \end{enumerate}
            
        \subsection{À Propos de / Au Sujet  de (about)}
            \begin{center}
                Il faut faire qqc à propos de la situation (We need to do something about the situation)
            \end{center}
        
        \subsection{Avec (with)}
            \begin{center}
                Il est avec son frère (He is with his brother)
            \end{center}
        \subsection{D'après (according to)}
            \begin{center}
                Il va pleuvoir d'après le météo (It is going to rain according to the weather forecast)
            \end{center}
        \subsection{Environ (approximately, about)}
            \begin{center}
                Il y a environ 10 personnes (There are about 10 people)
            \end{center}
        \subsection{Malgré (despite)}
            \begin{center}
                Il est sorti malgré la pluie (He went out despite the rain)
            \end{center}
        \subsection{Parmi (among)}
            \begin{center}
                Il est le meilleur parmi les trois (He is the best among the three)
            \end{center}
        \subsection{Sans (without)}
            \begin{center}
                Il est parti sans dire au revoir (He left without saying goodbye)
            \end{center}
        \subsection{Sauf (except)}
            \begin{center}
                Tout le monde est venu sauf lui (Everyone came except him)
            \end{center}
        \subsection{Selon (according to)}
            \begin{center}
                Selon le professeur, c'est une bonne idée (According to the professor, it is a good idea)
            \end{center}
        \subsection{Vers / Envers (towards, around + time)}
            \begin{center}
                Il est parti vers la maison (He left towards the house)
            \end{center}
    
    \section{Quand ne pas utiliser l'article}
        \begin{itemize}
            \item Profession: except after \textit{c'est}
                \begin{center}
                    Il est docteur \& C'est un docteur (He is a doctor)
                \end{center}

            \item After \textit{de}, if the noun is plural and after \underline{references of quantities}
                \begin{center}
                    Il y a beaucoup de gens (There are a lot of people)
                \end{center}

            \item After some prepositions, when talking in a general way
                \begin{itemize}
                    \item Avec:
                        \begin{center}
                            Elle cherche un appartement avec balcon (She is looking for an apartment with a balcony)
                        \end{center}

                    \item Sans:
                        \begin{center}
                            Tu ne te sentiras pas mieux sans étirements (You won't feel better without stretching)
                        \end{center}

                    \item Sous:
                        \begin{center}
                            L'Europe est sous eaupéenne (Europe is under European law)
                        \end{center}

                    \item Par:
                        \begin{center}
                            Il fait ça par habitude (He does it out of habit)
                        \end{center}

                    \item En: 
                        \begin{center}
                            Nous sommes en classe (We are in class)
                        \end{center}

                    \item After \textit{parler}, when talking about languages
                        \begin{center}
                            Elle parle anglais (She speaks English)
                        \end{center}
                \end{itemize}
        \end{itemize}
    
    \section{Verbs + Préposition}
        \begin{itemize}
            \item Aller (to go):
                \begin{itemize}
                    \item Aller vers + time (to go around + time)
                        \begin{center}
                            Il est allé vers 10 heures (He went around 10 o'clock)
                        \end{center}

                    \item Aller vers + location (to go toward + location)
                        \begin{center}
                            Il est allé vers la maison (He went toward the house)
                        \end{center}
                \end{itemize}
            
            \item Apprendre:
                \begin{itemize}
                    \item Apprendre à (to learn to):
                        \begin{center}
                            Il apprend à nager (He is learning to swim)
                        \end{center}
                    
                    \item Apprendre de (to learn from):
                        \begin{center}
                            Il apprend de ses erreurs (He learns from his mistakes)
                        \end{center}
                \end{itemize}
            
            \item Arriver:
                \begin{itemize}
                    \item Arriver à (to manage to / to succeed in):
                        \begin{center}
                            Il est arrivé à finir le travail (He managed to finish the work)
                        \end{center}
                    
                    \item Arriver de + location (to arrive from  + location):
                        \begin{center}
                            Il est arrivé de Paris (He arrived from Paris)
                        \end{center}
                \end{itemize}
            
            \item Chercher:
                \begin{itemize}
                    \item Chercher (to look for):
                        \begin{center}
                            Il cherche son chien (He is looking for his dog)
                        \end{center}
                    
                    \item Chercher à (to attempt to / to try to):
                        \begin{center}
                            Il cherche à comprendre (He is trying to understand)
                        \end{center}
                        
                    \item Chercher dans (to look in):
                        \begin{center}
                            Il cherche dans le sac (He is looking in the bag)
                        \end{center}
                \end{itemize}

            \item Commencer:
                \begin{itemize}
                    \item Commencer à (to begin to / to start to):
                        \begin{center}
                            Il commence à pleuvoir (It is starting to rain)
                        \end{center}
                    
                    \item Commencer par (to begin by / to start by):
                        \begin{center}
                            Il commence par un café (He starts with a coffee)
                        \end{center}
                \end{itemize}

            \item Compter (to expect / to intend):
                \begin{itemize}
                    \item Compter pour (to count / to be worth):
                        \begin{center}
                            Il compte pour deux (He counts for two)
                        \end{center}
                    
                    \item Compter sur (to count on):
                        \begin{center}
                            Il compte sur toi (He is counting on you)
                        \end{center}
                \end{itemize}
            
            \item Croire (to think / to believe):
                \begin{itemize}
                    \item Croire à (to believe something other than people):
                        \begin{center}
                            Il croit à la chance (He believes in luck)
                        \end{center}
                    
                    \item Croire en (to believe in), for god \& people:
                        \begin{center}
                            Il croit en Dieu (He believes in God)
                        \end{center}
                \end{itemize}
            
            \item Donner (to give):
                \begin{itemize}
                    \item Donner à (to give to):
                        \begin{center}
                            Il donne à manger au chien (He is feeding the dog)
                        \end{center}
                    
                    \item Donner de (to give of):
                        \begin{center}
                            Il donne de son temps (He gives of his time)
                        \end{center}
                    
                    \item Donner sur (to overlook / to open on):
                        \begin{center}
                            La fenêtre donne sur la rue (The window overlooks the street)
                        \end{center}
                \end{itemize}

            \item Être:
                \begin{itemize}
                    \item Être à (to belong to):
                        \begin{center}
                            Il est à moi (He is mine)
                        \end{center}
                    
                    \item Être pour (to be in favor of):
                        \begin{center}
                            Il est pour la paix (He is for peace)
                        \end{center}
                    
                    \item Être vers + time (to be around / near + time):
                        \begin{center}
                            Il est vers 10 heures (It is around 10 o'clock)
                        \end{center}
                \end{itemize}

            \item Finir:
                \begin{itemize}
                    \item Finir de (to finish):
                        \begin{center}
                            Il finit de manger (He finishes eating)
                        \end{center}
                        
                    \item Finir par (to end up):
                        \begin{center}
                            Il finit par accepter (He ends up accepting)
                        \end{center}
                \end{itemize}

            \item Parler:
                \begin{itemize}
                    \item Parler à (to leave): 
                        \begin{center}
                            Il parle à son ami (He is talking to his friend)
                        \end{center}    
                    
                    \item Partir dans + time (to leave in + time):
                        \begin{center}
                            Il part dans 10 minutes (He is leaving in 10 minutes)
                        \end{center}
                    
                    \item Partir pour (to leave for / to be off to):
                        \begin{center}
                            Il part pour Paris (He is leaving for Paris)
                        \end{center}
                \end{itemize}

            \item Tenir:
                \begin{itemize}
                    \item Tenir à (to insist / to be attached to):
                        \begin{center}
                            Il tient à son chien (He is attached to his dog)
                        \end{center}
                    
                    \item Tenir de (to take after):
                        \begin{center}
                            Il tient de son père (He takes after his father)
                        \end{center}
                \end{itemize}
            
            \item Venir:
                \begin{itemize}
                    \item Venir de (to have just):
                        \begin{center}
                            Il vient de manger (He has just eaten)
                        \end{center}
                    
                    \item Venir par (to come by):
                        \begin{center}
                            Il vient par la rue (He comes by the street)
                        \end{center}

                    \item Venir à (to come to):
                        \begin{center}
                            Il vient à la maison (He comes to the house)
                        \end{center}
                \end{itemize}
        \end{itemize}

\chapter{Questions}
    \section{Yes / No Questions}
    The following are ways to ask yes / no questions in French:
        \begin{enumerate}
            \item Raise your voice in the end (usually in spoken French):
                \begin{center}
                    Tu vas bien? (Are you okay?)
                \end{center}

            \item Add \textit{Est-ce que} in front of subject (kind of formal):
                \begin{center}
                    Est-ce que tu vas bien? (Are you okay?)
                \end{center}

            \item Inversion (formal):
                \begin{center}
                    Verb - subject / auxiliary - subject + past participle \\
                \end{center}
                For example:
                \begin{center}
                    Vas-tu bien? (Are you okay?) \\
                    Avez-vous vu le film? (Have you seen the movie?)
                \end{center}
                If we have two vowels, we add a \textit{t} in between:
                \begin{center}
                    Il aime le jus d'orange? (Does he like orange juice?) \\
                    Aime-t-il le jus d'orange? (Does he like orange juice?)
                \end{center}
        \end{enumerate}
    
    \section{Quand, Comment, Pourquoi, Où, Combien, Combien de}
        \begin{itemize}
            \item Formal: 
                \begin{center}
                    adverb + est-ce que / adverb + inversion
                \end{center}

            \item Informal: 
                \begin{center}
                    question + adverb / adverb + question
                \end{center}
            
            \item Examples:
                \begin{enumerate}
                    \item When are you leaving?
                        \begin{center}
                            \textbf{Formal.} Quand est-ce que tu pars? \\
                            \textbf{Informal.} Tu pars quand?
                        \end{center}
                    \item How are you?
                        \begin{center}
                            \textbf{Formal.} Comment est-ce que vous allez? / Comment allez-vous? \\
                            \textbf{Informal.} Comment vous allez? / Comment tu vas?
                        \end{center}
                    \item Why are they doing that?
                        \begin{center}
                            \textbf{Formal.} Pourquoi est-ce qu'ils font ça? / Pourpuoi font-ils ça? \\
                            \textbf{Informal.} Pourquoi ils font ça? 
                        \end{center}
                    \item Where do you live?
                        \begin{center}
                            \textbf{Formal.} Où est-ce que tu vis? / Où vis-tu? \\
                            \textbf{Informal.} Tu vis où?
                        \end{center}
                    \item How much did you receive?
                        \begin{center}
                            \textbf{Formal.} Combien est-ce que tu as reçu? / Combien as-tu reçu? \\
                            \textbf{Informal.} Tu as reçu combien? 
                        \end{center}
                    \item Add \textit{de / d'} after \textit{combien} when followed by a noun or adjective:
                        \begin{center}
                            How many friends do you have? \\
                            \textbf{Formal.} Combien d'amis est0ce que tu as? / Combien d'amis as-tu? \\
                            \textbf{Informal.} Tu as combien d'amis?
                        \end{center}
                \end{enumerate}
        \end{itemize}

    \section{Interrogative Pronouns}
        \begin{itemize}
            \item Qui
                \begin{enumerate}
                    \item As subject (who), no inversion:
                        \begin{center}
                            Qui \dots / Qui est-ce qui \dots
                        \end{center}
                        Example:
                            \begin{center}
                                Qui est-ce qui parle? \\
                                (Who is talking?)
                            \end{center}

                    \item As object (whom), inversion allowed:
                        \begin{center}
                            Qui \dots / Qui est-ce que \dots
                        \end{center}
                        Example:
                            \begin{center}
                                Qui est-ce que tu aimes? \\
                                (Who do you like?)
                            \end{center}
                \end{enumerate}
            
            \item Qui with proposition (whom): 
                \begin{enumerate}
                    \item À qui \dots / À qui est-ce que \dots
                        \begin{center}
                            À qui parles-tu? / à qui est-ce que tu parles? \\
                            (Who are you talking to?)
                        \end{center}

                    \item De qui \dots / De qui est-ce que \dots
                        \begin{center}
                            De qui est-ce que tu parles? \\
                            (Who are you talking about?)
                        \end{center}

                    \item Avec qui \dots / Avec qui est-ce que \dots
                        \begin{center}
                            Avec qui est-ce que tu vas? \\
                            (Who are you going with?)
                        \end{center}
                \end{enumerate}

            \item Que (what):
                \begin{enumerate}
                    \item As subject, no inversion: Qu'est-ce que \dots
                        \begin{center}
                            Qu'est-ce que tu fais? \\
                            (What are you doing?)
                        \end{center}

                    \item As object, inversion allowed, also can be changed to \textit{quoi} and put to the end:
                        \begin{center}
                            Que \dots / Qu'est-ce que \dots
                        \end{center}
                        Example:
                        \begin{center}
                            Que veux-tu manger? / Qu'est-ce que tu veux manger? \\
                            Tu veux manger quoi? \\
                            (What do you want to eat?)
                        \end{center}

                    \item Quoi = Que with proposition:
                        \begin{enumerate}
                            \item À quoi \dots / À quoi est-ce que \dots
                                \begin{center}
                                    À quoi penses-tu? / À quoi est-ce que tu penses? \\
                                    Tu pense à quoi? (informal) \\
                                    (What are you thinking about?)
                                \end{center}

                            \item De quoi \dots / De quoi est-ce que \dots / \dots de quoi? (informal):
                                \begin{center}
                                    De quoi parles-tu? / De quoi est-ce que tu parles? \\
                                    Tu parles de quoi? (informal) \\
                                    (What are you talking about?)
                                \end{center}

                            \item Avec quoi \dots / Avec quoi est-ce que \dots / \dots avec quoi? (informal):
                                \begin{center}
                                    Avec quoi est-ce que tu écris? / Avec quoi écris-tu? \\
                                    Tu écris avec quoi? (informal) \\
                                    (What are you writing with?)
                                \end{center}
                        \end{enumerate}
                \end{enumerate}

            \item Quel (which):
                \begin{center}
                    \begin{tabular}{|c|c|c|c|}
                        \hline
                        \textbf{Mas. Sin.} & \textbf{Fem. Sin.} & \textbf{Mas. Plu.} & \textbf{Fem. Plu.}\\
                        \hline
                        Quel & Quelle & Quels & Quelles \\
                        \hline
                    \end{tabular}
                \end{center}

                \begin{itemize}
                    \item Quel + noun (inversion \& \textit{est-ce que} allowed):
                        \begin{center}
                            \textcolor{red}{Quel} plat as-tu commandé? / \textcolor{red}{Quel} plat \textcolor{red}{est-ce que} tu as commandé? \\
                            Tu as commandé \textcolor{red}{quel} plat? (informal) \\
                            (Which dish did you order?)
                        \end{center}

                    \item Quel + verb (no inversion or \textit{est-ce que}):
                        \begin{center}
                            \textcolor{red}{Quel} est ton film préféré?
                        \end{center}

                    \item Quel with propositions:
                        \begin{enumerate}
                            \item De quel \dots / De quel est-ce que \dots
                                \begin{center}
                                    De quel pays est-il? / De quel est-ce que pays est-il? \\
                                    Il est de quel pays? (informal) \\
                                    (Which country is he from?)
                                \end{center}

                            \item À quel \dots / À quel est-ce que \dots
                                \begin{center}
                                    À quel hôtel allez-vous? / À quel est-ce que hôtel allez-vous? \\
                                    Vous allez à quel hôtel? (informal) \\
                                    (Which hotel are you going to?)
                                \end{center}

                            \item Avec quel \dots / Avec quel est-ce que \dots
                                \begin{center}
                                    Avec quel crayon écrivez-vous? / Avec quel est-ce que crayon écrivez-vous? \\
                                    Vous écrivez avec quel crayon? (informal) \\
                                    (Which pencil are you writing with?)
                                \end{center}
                        \end{enumerate}

                    \item Case for follow up question, quel + noun:
                        \begin{center}
                            Qu'est-ce que tu penses de mon idée? \\
                            \textcolor{red}{Quelle} idée? \\
                        \end{center}
                \end{itemize}

            \item Lequel (which one):
                \begin{itemize}
                    \item Contraction for \textit{Le pull + quel = le + quel = Lequel}
                    \begin{center}
                        \begin{tabular}{|c|c|c|c|c|}
                            \hline
                            & \textbf{Mas. Sin.} & \textbf{Fem. Sin.} & \textbf{Mas. Plu.} & \textbf{Fem. Plu.}\\
                            \hline
                            $\emptyset$ & Lequel & Laquelle & Lesquels & Lesquelles \\
                            \hline
                            à + lequel & Auquel & À laquelle & Auxquels & Auxquelles \\
                            \hline
                            de + lequel & Duquel & De laquelle & Desquels & Desquelles \\
                            \hline
                        \end{tabular}
                    \end{center}
                \end{itemize}
        \end{itemize}

    \section{Tag Questions}
        Aren't you? Isn't she? \dots
        \begin{itemize}
            \item Use \textit{N'est-ce que}:
                \begin{center}
                    Il est bien arrivé, n'est-ce pas? \\
                    (He arrived well, didn't he?)
                \end{center}

            \item For more critical way, use \textit{Hein? / Non?}:
                \begin{center}
                    Il est bien arrivé, hein? \\
                    C'est ta soeur, non? \\
                    (That's your sister, right?)
                \end{center}
            
            \item D'accord (ok):
                \begin{center}
                    On part bientôt, d'accord? \\
                    (We are leaving soon, ok?) \\
                    Tu vas te coucher dans 5 minutes, d'accord? \\
                    (You are going to bed in 5 minutes, ok?)
                \end{center}

            \item Si? (use in negative questions):
                \begin{center}
                    Ils ne sont pas mariés, si? \\
                    (They are not married, are they?)
                \end{center}
        \end{itemize}

    \section{Answer a Negative Quesiton}
        Use \textit{si} instead of \textit{oui}:
            \begin{center}
                Ils ne sont pas mariés? \\
                Si, ils sont mariés depuis quelques années. 
            \end{center}   
            
\chapter{Les Adjectifs}
    \begin{itemize}
        \item Pronounciation changes for consonant endings.
        \item Regular endings:
            \begin{center}
                \begin{tabular}{|c|c|c|c|c|}
                    \hline
                    & \textbf{M.S.} & \textbf{F.S.} & \textbf{M.P.} & \textbf{F.P.} \\
                    \hline
                    \textbf{-E} & $\emptyset$ & $\emptyset$ & -s & -s \\
                    \hline
                    \textbf{-É, -I, -U} & $\emptyset$ & -e & -s & -es \\
                    \hline
                    \textbf{-T, -D, -N, -S} & $\emptyset$ & -e & -s / $\emptyset$ & -es \\
                    \hline
                    \textbf{-EUX} & $\emptyset$ & -x $\to$ -se & $\emptyset$ & -x $\to$ -ses \\
                    \hline
                    \textbf{-F} & $\emptyset$ & -f $\to$ -ve & -s & -s $\to$ -ves \\
                    \hline
                    \textbf{-ER} & $\emptyset$ & -er $\to$ -ère & -s & -er $\to$ -ères \\
                    \hline
                    \textbf{-AL} & $\emptyset$ & -e & -al $\to$ -aux & -es \\
                    \hline
                \end{tabular}
            \end{center}
        \item Last consonant double: when simply adding -e results in another word. Let $*$ denote the consonant, we have
            \begin{center}
                \begin{tabular}{|c|c|c|c|}
                    \hline 
                    \textbf{M.S.} & \textbf{F.S.} & \textbf{M.P.} & \textbf{F.P.} \\
                    \hline
                    $\emptyset$ & -*e & -s / $\emptyset$ & -*es \\
                    \hline
                \end{tabular}
            \end{center}
        \item Irrèguliers:
            \begin{center}
                \begin{tabular}{|c|c|c|c|c|}
                    \hline
                    \textbf{M.S.} & \textbf{F.S.} & \textbf{M.P.} & \textbf{F.P.} & \textbf{English} \\
                    \hline
                    Blanc & Blanche & Blancs & Blanches & White \\
                    \hline
                    Dou\textcolor{red}{x} & Dou\textcolor{red}{ce} & Dou\textcolor{red}{x} & Dou\textcolor{red}{ces} & Soft \\
                    \hline
                    Favori & Favori\textcolor{red}{te} & Favori\textcolor{red}{s} & Favori\textcolor{red}{tes} & Favorite \\
                    \hline
                    Fou & Folle & Fous & Folles & Crazy \\
                    \hline
                    Frais & Fra\textcolor{red}{îche} & Frai\textcolor{red}{s} & Fra\textcolor{red}{îches} & Fresh \\
                    \hline
                    Long & Long\textcolor{red}{ue} & Longu\textcolor{red}{s} & Long\textcolor{red}{ues} & Long \\
                    \hline
                    Sec & S\textcolor{red}{èche} & Sec\textcolor{red}{s} & S\textcolor{red}{èches} & Dry \\
                    \hline
                    Inquiet & Inqui\textcolor{red}{ète} & Inqui\textcolor{red}{ets} & Inqui\textcolor{red}{ètes} & Worried \\
                    \hline
                \end{tabular}
            \end{center}
        \item Adjectives with a $5^{th}$ choice (if word after starts with a vowel or silent h, for M.S. only):
            \begin{center}
                \begin{tabular}{|c|c|c|c|c|}
                    \hline
                    \textbf{M.S.} & \textbf{F.S.} & \textbf{M.P.} & \textbf{F.P.} & \textbf{English} \\
                    \hline
                    nouveau / nouvel & nouvelle & nouveaux & nouvelles & New \\
                    \hline
                    Vieux / Vieil & Vieille & Vieux & Vieilles & Old \\
                    \hline
                    beau / bel & belle & beaux & belles & Beautiful \\
                    \hline
                \end{tabular}
            \end{center}
        \item Tout (all / whole): 
            \begin{itemize}
                \item Tout, toute, tous, toutes
                \item Usually placed after the noun, for example:
                    \begin{center}
                        Tous les jours (Every day) \\
                        Toute la journée (All day) \\
                        Tous les deux (Both of them) \\
                        Toutes les deux (Both of them)
                    \end{center}
            \end{itemize}
        \item Adjectives placed before noun:
            \begin{itemize}
                \item BANGS (beauty, age, number, goodness, size).
                    \begin{enumerate}
                        \item Beauty: beau, joli \dots
                        \item Age: Heune, vieux, nouveau \dots
                        \item Number: All numbers, premier, second \dots, dernier
                        \item Goodness: Bon, Mauvais, Meilleur, Gentil \dots
                        \item Size: Petit, Grand, Gros, Long \dots
                    \end{enumerate}
            \end{itemize}
        \item More than one adjective:
            \begin{itemize}
                \item When both adjectives should be placed after the noun, add \textit{et} in between.
            \end{itemize}
        \item Adjectives with 2 different meanings:
            \begin{itemize}
                \item Before the noun implies figurative Meaning.
                \item After the noun implies literal meaning.
                \begin{enumerate}
                    \item Ancien / Ancienne:
                        \begin{center}
                            Mon \textcolor{red}{ancienne} voiture (My old car) \\
                            Une voiture \textcolor{red}{ancienne} (An antique car)
                        \end{center}
                    \item Dernier / Dernière:
                        \begin{center}
                            Le \textcolor{red}{dernier} mardi (The last Tuesday) \\
                            Mardi \textcolor{red}{dernier} (last Tuesday)
                        \end{center}
                    \item Drôle:
                        \begin{center}
                            Une \textcolor{red}{drôle} de personne (A strange person) \\
                            Une personne \textcolor{red}{drôle} (A funny person)
                        \end{center}
                    \item Pauvre:
                        \begin{center}
                            Un \textcolor{red}{pauvre} enfant (A poor child (pity)) \\
                            Un enfant \textcolor{red}{pauvre} (A poor child (money))
                        \end{center}
                    \item Propre:
                        \begin{center}
                            Ma \textcolor{red}{propre} maison (My own house) \\
                            Une assiette \textcolor{red}{propre} (A clean plate)
                        \end{center}
                \end{enumerate}
            \end{itemize}
    \end{itemize}

\chapter{Comparatives}
    Used for comparing two things.
    \begin{itemize}
        \item Aussi + adjective + que (as \dots as \dots):
            \begin{center}
                Il est aussi grand que moi (He is as tall as me)
            \end{center}
        \item Plus + adjective + que (more \dots than \dots):
            \begin{center}
                Il est plus grand que moi (He is taller than me)
            \end{center}
        \item Mois + adjective + que (Less \dots than \dots)
            \begin{center}
                Il est moins grand que moi (He is less tall than me)
            \end{center}
        \item Autant que (as much as / as many as)
            \begin{center}
                Il a autant d'argent que moi (He has as much money as me)
            \end{center}
    \end{itemize}
    If we know the subject of the comparison, \textit{que} \dots can be omitted.

    \section{Comparative + noun}
        \begin{itemize}
            \item Autant de + noun + que (as much \dots as)
                \begin{center}
                    Il a autant d'argent que moi (He has as much money as me)
                \end{center}
            \item Plus de + noun + que (more \dots than)
                \begin{center}
                    Il a plus d'amis que moi (He has more friends than me)
                \end{center}
            \item Moins de + noun + que (less \dots than)
                \begin{center}
                    Il a moins de temps que moi (He has less time than me)
                \end{center}
        \end{itemize}

    \section{Superlative}
        Talk about extreme, do not use for comparing two things.
        \begin{itemize}
            \item le / la / les plus + adjective (the most)
                \begin{center}
                    C'est le plus grand (He is the tallest)
                \end{center}
            \item le / la / les moins + adjective (the least)
                \begin{center}
                    C'est le moins cher (It is the least expensive)
                \end{center}
        \end{itemize}
        To add context: add \textit{de} at the end of the adjective:
        \begin{center}
            C'est le plus grand de la classe (He is the tallest in the class)
        \end{center}

    \section{Irregular Comparatives}
        \begin{center}
            \begin{tabular}{|c|c|c|}
                \hline
                \textbf{Irregular} & \textbf{Comparative} & \textbf{Superlative} \\
                \hline
                Bon & Meilleur & Le meilleur \\
                \hline
                Mauvais & Pire & Le pire \\
                \hline
                Petit & Moindre & Le moindre \\
                \hline
            \end{tabular}
        \end{center}

\chapter{Les Adj. Indéfinis}
    Help to describe a noun in a non-specific way, always placed before the noun.
    \begin{center}
        \begin{tabular}{|c|c|c|c|c|}
            \hline
            \textbf{M.S.} & \textbf{F.S.} & \textbf{M.P.} & \textbf{F.P.} & \textbf{English} \\
            \hline
            Autre & Autre & Autres & Autres & Other \\
            \hline
            Certain & Certaine & Certains & Certaines & Certain \\
            \hline
            Chaque & Chaque & Chaque & & \\
            \hline
            & & Divers & Diverses & Various \\
            \hline
            & & Plusieurs & Plusieurs & Several \\
            \hline
            Quelque & Quelque & Quelques & Quelques & Some \\
            \hline
            Tel & Telle & Tels & Telles & Such \\
            \hline
            Tout & Toute & Tous & Toutes & All \\
            \hline
        \end{tabular}
    \end{center}

\chapter{The Direct Object}
    \begin{itemize}
        \item A noun or a group of words including a noun.
        \item Referring to people or things, usually placed directly after the verb.
        \item Found by asking the direct questions \textit{Quoi} or \textit{Qui}.
        \item The format is:
            \begin{center}
                Subject + verb + direct object
            \end{center}
    \end{itemize}
    \begin{enumerate}
        \item The direct object + Quoi:
            \begin{center}
                Je lis quoi? (What am I reading?) \\
                Un livre
            \end{center}
            \textit{Un livre} is the direct object.

        \item The direct object + Qui:
            \begin{center}
                Il appelle qui? (Who is he calling?) \\
                Ashley
            \end{center}
            \textit{Ashley} is the direct object.
    \end{enumerate}

    \section{The Direct Object Pronoun}
        \begin{itemize}
            \item Replace the direct object, usually placed before the verb.
            \item The format is:
                \begin{center}
                    Subject + direct object pronoun + verb
                \end{center}
        \end{itemize}
        \begin{enumerate}
            \item Singular: me / m', te / t', le / l', la / l'
            \item Plural: nous, vous, les
            \item Example:
                \begin{center}
                    Elle lave sa nouvelle voiture $\to$ Elle la lave \\
                    (She washes her new car $\to$ She washes it)
                \end{center}
        \end{enumerate}

        \subsection{REDCAP}
            6 verbs that don't take a preposition in French and are followed by a direct object:
            \begin{enumerate}
                \item Regarder qqch / qqn (to look at):
                    \begin{center}
                        Je regarde le film (I am watching the movie)
                    \end{center}
                \item Écouter qqch / qqn (to listen to):
                    \begin{center}
                        J'écoute la musique (I am listening to music)
                    \end{center}
                \item Demander qqch / qqn (to ask for):
                    \begin{center}
                        Je demande le menu (I ask for the menu)
                    \end{center}
                \item Chercher qqch / qqn (to look for):
                    \begin{center}
                        Il cherche son chien (He is looking for his dog)
                    \end{center}
                \item Attendre qqch / qqn (to wait for):
                    \begin{center}
                        J'attends le bus (I am waiting for the bus)
                    \end{center}
                \item Payer qqch (to pay for):
                    \begin{center}
                        Je paie le repas (I pay for the meal)
                    \end{center}
            \end{enumerate}

    \section{Direct Object Pronoun \& Negation}
        The format is:
            \begin{center}
                subject + ne + direct object pronoun + verb + pas
            \end{center}
        Examples:
            \begin{center}
                Est-ce que tu regardes \textcolor{red}{la télévision}? \\
                Non, je ne \textcolor{red}{la} regarde pas. 
            \end{center}

    \section{Direct Object Pronoun \& Passé Composé}
        \begin{enumerate}
            \item DOP + avoir + past participle + e / s / es:
                \begin{center}
                    J'ai commandé \textcolor{red}{un thé} $\to$ Je \textcolor{red}{l'}ai commandé \\
                    (I ordered a tea $\to$ I ordered it) \\
                    J'ai commandé \textcolor{red}{deux thé} $\to$ Je \textcolor{red}{les} ai commandées \\
                    (I ordered two teas $\to$ I ordered them)
                \end{center}
            
            \item A relative pronoun: \textit{Que / Qu'}:
                \begin{center}
                    J'ai commandé \textcolor{red}{un thé} $\to$ Le thé \textcolor{red}{que} j'ai commandé \\
                    (I ordered a tea $\to$ The tea that I ordered) \\
                    J'ai commandé \textcolor{red}{deux glaces} $\to$ Les deux glaces \textcolor{red}{que} j'ai commandé\textcolor{red}{es} \\
                    (I ordered two ice creams $\to$ The two ice creams that I ordered)
                \end{center}
                The reason for adding \textit{es} is because the noun is feminine and plural.
        \end{enumerate}

\chapter{The Indirect Object}
    \begin{itemize}
        \item Usually preceded by the preposition \textit{à} or \textit{pour}.
        \item Found by asking the question \textit{À qui} or \textit{Pour qui}.
        \item Placed before the verb.
        \item The format is:  
            \begin{center}
                subject + IOP + verb
            \end{center}
        \item Examples:
            \begin{center}
                Je donne un conseil à qui? (To whom am I giving advice?) \\
                À Jean (To Jean) \\
                Elle achète un cadeau à qui? (To whom is she buying a gift?) \\
                À ses parents (To her parents) \\
                L'infirmier donne les médicaments \textcolor{red}{au patient} $\to$ L'infirmier \textcolor{red}{lui} donne les médicaments \\
                (The nurse gives the medicine to the patient $\to$ The nurse gives him the medicine)
            \end{center}
    \end{itemize}
    \begin{enumerate}
        \item Singular: me / m', te / t', lui
        \item Plural: nous, vous, leur
    \end{enumerate}
    Verbs that has \textit{à} as a preposition:
        \begin{center}
            \begin{tabular}{|c|c|}
                \hline
                Conseiller à qqn & To advise someone \\
                Dire à qqn & To tell someone \\
                Demander à qqn & To ask someone \\
                Obéir à qqn & To obey someone \\
                Rendre visite à qqn & To visit someone \\
                Répondre à qqn & To answer someone \\
                \dots & \dots \\
                \hline
            \end{tabular}
        \end{center}
    There 3 verbs are followed by an emphatic pronoun:
        \begin{center}
            \begin{tabular}{|c|c|}
                \hline
                Être à & To belong to \\
                Penser à & To think about \\
                Faire attention à & To pay attention to \\
                \hline
            \end{tabular}
        \end{center}
    Examples:
        \begin{center}
            \begin{tabular}{|c|c|}
                \hline
                C'est à moi & It is mine \\
                Je pense à toi & I am thinking about you \\
                Fais attention à toi & Pay attention to yourself \\
                \hline
            \end{tabular}
        \end{center}

    \section{Negation}
        The format is:
            \begin{center}
                subject + ne + IOP + verb + pas
            \end{center}
        Examples:
            \begin{center}
                Est-ce que tu offres souvent des fluers \textcolor{red}{à ta femme}? \\
                Non, je ne \textcolor{red}{lui} offre jamais de fleurs \\
                (Do you often offer flowers to your wife? No, I never offer her flowers)
            \end{center}

    \section{The Pronoun Y}
        The format is:
            \begin{center}
                subject + y + verb
            \end{center}
        \begin{enumerate}
            \item For places, \textit{Y} will be used to replace the following prepositions:
                \begin{center}
                    à, en, dans, chez \\
                    Ex. On est \textcolor{red}{au magasin} $\to$ On \textcolor{red}{y} est \\
                    (We are at the store $\to$ We are there)
                \end{center}
            \item For things, replace:
                \begin{center}
                    à / au / aux / à l' / à la \\
                    Ex. Je pense \textcolor{red}{à mes vacances} $\to$ J'\textcolor{red}{y} pense \\
                    (I am thinking about my vacation $\to$ I am thinking about it) \\
                    Ex. Je vais réfléchir \textcolor{red}{à cela} $\to$ Je vais \textcolor{red}{y} réfléchir \\
                    (I am going to think about that $\to$ I am going to think about it)
                \end{center}
            \item Expressions: \textit{Il y a}, \textit{vas-y / allons-y} (let's go), \textit{viens-y / venez-y} (come here)
                \begin{center}
                    Il y a un problème (There is a problem) \\
                    Vas-y (Go there) \\
                    Allons-y (Let's go there) \\
                    Viens-y (Come here)
                \end{center}
        \end{enumerate}

    \section{The Pronoun EN}
        An object including the preposition \textit{de}, and expressions of quantity, which are most likely to include a partitive article: \textit{de, du, de la, des} \dots
        \begin{center}
            subject + en + verb
        \end{center}
        Examples:
        \begin{center}
            Est-ce que tu as \textcolor{red}{du lait}? \\
            Oui, j'\textcolor{red}{en} ai \\
            (Do you have milk? Yes, I have some) \\
            Tu as beaucoup \textcolor{red}{d'argent} $\to$ Tu \textcolor{red}{en} as beaucoup \\
            (You have a lot of money $\to$ You have a lot) \\
            J'ai \textcolor{red}{deux verres} $\to$ J'\textcolor{red}{en} ai deux \\
            (I have two glasses $\to$ I have two)
        \end{center}
        For the last example, keep the number.
        
        \begin{enumerate}
            \item A few verbs taking EN:
                \begin{enumerate}
                    \item Avoir besoin de (to need):
                        \begin{center}
                            J'ai besoin de \textcolor{red}{conseils} $\to$ J'\textcolor{red}{en} ai besoin \\
                            (I need advice $\to$ I need some)
                        \end{center}
                    \item Avoir envie de (to want):
                        \begin{center}
                            J'ai envie de \textcolor{red}{frites} $\to$ J'\textcolor{red}{en} ai envie \\
                            (I want fries $\to$ I want some)
                        \end{center}
                    \item S'occuper de (to take care of):
                        \begin{center}
                            Je m'occupe de \textcolor{red}{mes enfants} $\to$ Je m'\textcolor{red}{en} occupe \\
                            (I take care of my children $\to$ I take care of them)
                        \end{center}
                    \item Parler de (to talk about):
                        \begin{center}
                            Je parle de \textcolor{red}{mon travail} $\to$ Je \textcolor{red}{en} parle \\
                            (I talk about my work $\to$ I talk about it)
                        \end{center}
                \end{enumerate}

            \item A few expressions with EN:
                \begin{enumerate}
                    \item J'en ai marre / J'en ai assez (I am fed up):
                        \begin{center}
                            J'en ai marre de \textcolor{red}{cette situation} \\
                            (I am fed up with this situation)
                        \end{center}
                    \item Je m'en moque / Je m'en fous (I don't care)
                    \item Je m'en vais (I am leaving)
                    \item Je suis en route (I am on my way)
                    \item Ne t'en fais pas (Don't worry)
                    \item Je suis en train de \dots (I am in the process of \dots)
                    \item Qu'est-ce que tu en passes? (What do you think?)
                \end{enumerate}
        \end{enumerate}

\chapter{Order of Using IOP \& DOP}
    The rule is me first, object second, other people third, and \textit{y} and \textit{en} last.
    \begin{center}
        \begin{tabular}{|c|c|}
            \hline
            \textbf{1st} & me / m', te / t', nous, vous \\
            \hline
            \textbf{2nd} & le / l', la / l', les \\
            \hline
            \textbf{3rd} & lui, leur \\
            \hline
            \textbf{4th} & y \\
            \hline
            \textbf{5th} & en \\
            \hline
        \end{tabular}
    \end{center}
    Examples:
    \begin{enumerate}
        \item Elle lit une histoire à \dots 
            \begin{tabular}{|c|c|}
                \hline
                Elle \textcolor{red}{me la} lit & She reads it to me \\
                Elle \textcolor{red}{te la} lit & She reads it to you \\
                Elle \textcolor{red}{nous la} lit & She reads it to us \\
                Elle \textcolor{red}{vous la} lit & She reads it to you \\
                Elle \textcolor{red}{leur la} lit & She reads it to them \\
                Elle \textcolor{red}{lui la} lit & She reads it to him / her \\
                \hline
            \end{tabular}

        \item Je donne des oranges à \dots 
            \begin{tabular}{|c|c|}
                \hline
                Je \textcolor{red}{t'en} donne & I give you some \\
                Je \textcolor{red}{vous en} donne & I give you some \\
                Je \textcolor{red}{lui en} donne & I give him / her some \\
                Je \textcolor{red}{leur en} donne & I give them some \\
            \end{tabular}
        
        \item Je regions \dots à Londres 
            \begin{tabular}{|c|c|}
                \hline
                Je \textcolor{red}{t'y} regoins & I am going there \\
                Je \textcolor{red}{vous y} regoins & I am going there \\
                Le / La: Je \textcolor{red}{l'y} regoins & I am going there \\
                Je \textcolor{red}{les y} regoins & I am going there \\
                \hline
            \end{tabular}
    \end{enumerate}

\chapter{Impératif \& Pronouns}
    \begin{enumerate}
        \item Verb pronoun: placed after the verb, \textit{me $\to$ moi, te $\to$ toi}, link with hyphen.
            \begin{center}
                Regarde-moi quand je te parle (Look at me when I talk to you) \\
                Tais-toi (Shut up) \\
                Allons-y (Let's go) \\
                Prends-en (Take some)
            \end{center}
        \item Negation:
            \begin{center}
                Ne + pronoun + verb + pas \\
                Ne t'inquiète pas pour ça (Don't worry about that) \\
                Ne me compte pas dans la liste (Don't count me in the list)
            \end{center}

    \end{enumerate}

    \section{The Emphatic Pronouns}
    Use to emphasize, to put the accent on a person. Only use for people. We have 
        \begin{center}
            moi, toi, lui, elle, nous, vous, $\underbrace{\text{eux, elles,}}_{\text{m / f for them}}$ soi (oneself)
        \end{center}
    Example ues cases:
        \begin{enumerate}
            \item On their own or with an adverb:
                \begin{center}
                    Qui a les meilleurs points? \textcolor{red}{Moi}. \quad (Who has the best points? Me) \\
                    J'ai faim! \textcolor{red}{Moi} aussi. \quad (I am hungry! Me too) 
                \end{center}
            \item After preposition such as \textit{avec, pour, sans, chez}:
                \begin{center}
                    Est-ce qu'il est avec \textcolor{red}{toi}? \quad (Is he with you?) \\
                    Je ne sais pas clormir sans \textcolor{red}{lui}. \quad (I can't sleep without him) \\
                    C'est pour \textcolor{red}{elle}? \quad (Is it for her?) \\
                    On est chez \textcolor{red}{moi}. \quad (We are at my place)
                \end{center}
            \item After \textit{à, de}:
                \begin{center}
                    On n'a pas besoin de \textcolor{red}{toi}. \quad (We don't need you) 
                \end{center}
            \item With comparisons:
                \begin{center}
                    Elle est plus âgée que \textcolor{red}{nous}. \quad (She is older than us)
                \end{center}
            \item With \textit{ne \dots que}:
                \begin{center}
                    Il ne pense qu'à \textcolor{red}{lui}! \quad (He only thinks about himself) \\
                    Je n'ai vu qu'\textcolor{red}{elle}. \quad (I only saw her)
                \end{center}
                The preposition \textit{à} and \textit{de} might be after \textit{que}, depending on the verb.
            \item To emphasize:
                \begin{center}
                    \textcolor{red}{Toi}, tu es grand, mais \textcolor{red}{elle}, elle est petite. \\
                    (You are tall, but she is short)
                \end{center}
            \item To repeat the subject:
                \begin{center}
                    Je n'ai pas envie d'y aller, \textcolor{red}{moi}. \quad (I don't want to go, me)
                \end{center}
            \item When the subject is made out of a pronoun and a noun, or two pronouns:
                \begin{center}
                    Mes parents et \textcolor{red}{moi} habitons ensemble. \\
                    (My parents and I live together)
                \end{center}
            \item After \textit{c'est, ce sont}:
                \begin{center}
                    C'est \textcolor{red}{moi} qui ai gagné. \quad (It is me who won) \\
                    Ce sont \textcolor{red}{eux} qui ont perdu. \quad (It is them who lost)
                \end{center}
        \end{enumerate}

        \subsection{How to use \textit{Soi}}
            Use \textit{soi} with words such as \textit{tout le monde, chacun, personne, on}, or impersonal sentences.
            \begin{enumerate}
                \item After \textit{on}:
                    \begin{center}
                        On doit penser à \textcolor{red}{soi}. \quad (One must think about oneself)
                    \end{center}
                \item After \textit{tout le monde}:
                    \begin{center}
                        Tout le monde doit s'occuper de \textcolor{red}{soi}. \\ 
                        (Everyone must take care of themselves)
                    \end{center}
                \item After \textit{chacun}:
                    \begin{center}
                        Chacun doit penser à \textcolor{red}{soi}. \quad (Everyone must think about themselves)
                    \end{center}
                \item After \textit{personne}:
                    \begin{center}
                        Personne ne s'occupe de \textcolor{red}{soi}. \quad (No one takes care of themselves)
                    \end{center}
                \item Impersonal sentences:
                    \begin{center}
                        Il faut penser à \textcolor{red}{soi}. \quad (One must think about oneself)
                    \end{center}
            \end{enumerate}
        
        \subsection{\textit{Même} or \textit{Mêmes}}
            Means \textit{same} or \textit{self}.
            \begin{center}
                \begin{tabular}{|c|c|c|c|c|}
                    \hline
                    moi-même & toi-même & lui-même & elle-même & soi-même \\
                    \hline
                    nous-mêmes & vous-mêmes & eux-mêmes & elles-mêmes & \\
                    \hline
                \end{tabular}
            \end{center}
            Example:
            \begin{center}
                C'est important de faire les choses \textcolor{red}{soi-même}. \\
                (It is important to do things oneself)
            \end{center}
    
    \section{The Possessive Pronouns}
        possessive adjective + noun = possessive pronoun
        \begin{center}
            \begin{tabular}{|c|c|c|c|c|}
                \hline
                & \textbf{M.S.}(le + \rule{0.5cm}{0.1mm}) & \textbf{F.S.}(la + \rule{0.5cm}{0.1mm}) & \textbf{M.P.}(les + \rule{0.5cm}{0.1mm}) & \textbf{F.P.}(les + \rule{0.5cm}{0.1mm}) \\
                \hline
                mine: & mien & mienne & miens & miennes \\
                \hline
                yours: & tien & tienne & tiens & tiennes \\
                \hline
                his / hers / its: & sien & sienne & siens & siennes \\
                \hline
                ours: & nôtre & nôtre & nôtres & nôtres \\
                \hline
                yours: & vôtre & vôtre & vôtres & vôtres \\
                \hline
                theirs: & leur & leur & leurs & leurs \\
                \hline
            \end{tabular}
        \end{center}

    \section{Demonstrative Pronoun}
        Demonstrative adjective + noun = demonstrative pronoun / ci / là 
        \begin{center}
            \begin{tabular}{|c|c|c|c|c|}
                \hline
                & \textbf{M.S.}: celui & \textbf{F.S.}: celle & \textbf{M.P.}: ceux & \textbf{F.P.}: celles \\
                \hline
                This one: & celui-ci & celle-ci & ceux-ci & celles-ci \\
                \hline
                That one: & celui-là & celle-là & ceux-là & celles-là \\
                \hline
            \end{tabular}
        \end{center}
        For example:
        \begin{center}
            Lequel est-ce que tu veux? Celui-ci. \quad (Which one do you want? This one) \\
            Ceux-ci sont magnifiques! \quad (These are magnificent)
        \end{center}

        \begin{itemize}
            \item If have both -ci \& -là, then -ci first than -là. -ci \& -là can also be added to nouns to express a stronger opinion.
                \begin{center}
                    Ces couleurs\textcolor{red}{-ci} sont plus vives que \textcolor{red}{celles-là}. \\
                    (These colors are brighter than those) \\
                    J'aurais acheté \textcolor{red}{celui-ci} mais je préfère \textcolor{red}{celui-là}. \\
                    (I would have bought this one but I prefer that one)
                \end{center}

            \item Followed by a relative pronoun, translates to \textit{the one who, the one that, those who, those that}:
                \begin{center}
                    \textcolor{red}{Celui} qui est là est mon frère. \quad (The one who is there is my brother) \\
                    C'est \textcolor{red}{celui} que je t'ai donné. (It's the one I gave you)
                \end{center}
        \end{itemize}

    \section{The Indefinite Demonstrative Pronouns}
        \begin{center}
            \begin{tabular}{|c|c|}
                \hline
                ce & it, this, that, there, those \\
                \hline
                ceci & this \\
                \hline
                cela & that \\
                \hline
                ça & this, that \\
                \hline
            \end{tabular}
        \end{center}

        \begin{itemize}
            \item \textit{ce} is commonly found with the verb \textit{être} in expressions \textit{c'est} or \textit{ce sont}.
                \begin{center}
                    \textcolor{red}{C'est} un bon film. \quad (It is a good movie) \\
                    \textcolor{red}{Ce sont} des bons films. \quad (They are good movies)
                \end{center}   
            \item \textit{ceci} and \textit{cela} with other verbs than \textit{être}.
                \begin{center}
                    \textcolor{red}{Cela} devrait être possible \quad (That should be possible) \\
                    \textcolor{red}{Ceci} derait fonctionner mieux \quad (This should work better)
                \end{center}
            \item \textit{ceci} and \textit{cela} can be direct object.
                \begin{center}
                    As-tu vu \textcolor{red}{ceci}? \quad (Have you seen this?) \\
                \end{center}
            \item \textit{ça} is abbreviation for \textit{ceci} and \textit{cela}, used for speaking. Common expressions include:
                \begin{center}
                    Ça va? \quad (How are you?) \\
                    Ça me plaît \quad (I like it) \\
                    Ça y est \quad (That's it/done/finished) \\
                    Ça suffit \quad (That's enough) \\
                    Ça craint \quad (That sucks) \\
                    Ça ne me regarde pas \quad (That's none of my business)
                \end{center}
        \end{itemize}

    \section{The Indefinite Pronouns}
        \begin{enumerate}
            \item Un autre (another one):
                \begin{center}
                    \begin{tabular}{|c|c|}
                        \hline
                        \textbf{M.S.} & \textbf{F.S.} \\
                        \hline
                        un autre & une autre \\
                        \hline
                    \end{tabular}
                \end{center}
                \begin{center}
                    Ce café est délicieux, je vais en prendre un \textcolor{red}{un autre}. \\ 
                    (This coffee is delicious, I will take another one)
                \end{center}

            \item D'autre (others):
                \begin{center}
                    \begin{tabular}{|c|c|}
                        \hline
                        \textbf{M.P.} & \textbf{F.P.} \\
                        \hline
                        d'autres & d'autres \\
                        \hline
                    \end{tabular}
                \end{center}
                \begin{center}
                    Beaucoup d'employés arrivent à l'heure, \textcolor{red}{d'autres} arrivent en retard. \\
                    (Many employees arrive on time, others arrive late)
                \end{center}

            \item Certain (certain ones): 
                \begin{center}
                    \begin{tabular}{|c|c|}
                        \hline
                        \textbf{M.P.} & \textbf{F.P.} \\
                        \hline
                        certains & certains \\
                        \hline
                    \end{tabular}
                \end{center}
                \begin{center}
                    Ces peintures sont incroyables mais \textcolor{red}{certaines} sont un peu bizarres. \\
                    (These paintings are increadible but certain ones are a little bit weird)
                \end{center}

            \item Chacun (each one): 
                \begin{center}
                    \begin{tabular}{|c|c|}
                        \hline
                        \textbf{M.S.} & \textbf{F.S.} \\
                        \hline
                        chacun & chacune \\
                        \hline
                    \end{tabular}
                \end{center}
                \begin{center}
                    \textcolor{red}{Chacun} prépare sa valise. \\
                    (Each one prepares his suitcase)
                \end{center}

            \item Plusieurs (several):
                \begin{center}
                    \begin{tabular}{|c|c|}
                        \hline
                        \textbf{M.P.} & \textbf{F.P.} \\
                        \hline
                        plusieurs & plusieurs \\
                        \hline
                    \end{tabular}
                \end{center}
                \begin{center}
                    Les habitations ont été inondées. \textcolor{red}{Plusiers} sont inhabitables. \\
                    (The houses have been flooded. Several are uninhabitable)
                \end{center}

            \item Quelque chose, only masculine singular:
                \begin{center}
                    \textcolor{red}{Quelque chose} me dit qu'il sera en tard. \\
                    (Something tells that me he will be late)
                \end{center}

            \item Quelqu'un (someone), only masculine singular:
                \begin{center}
                    \textcolor{red}{Quelqu'un} est venu à la porte. \\
                    (Someone came to the door)
                \end{center}
            \item Quelques-uns (some/a few):
                \begin{center}
                    \begin{tabular}{|c|c|}
                        \hline
                        \textbf{M.P.} & \textbf{F.P.} \\
                        \hline
                        quelques-uns & quelques-unes \\
                        \hline
                    \end{tabular}
                \end{center}
                \begin{center}
                    J'ai téléphoné à mes amis, quelques-uns ont répondu. \\
                    (I called my friends, some answered)
                \end{center}
            \item Quiconque (whoever/anyone), only masculine singular:
                \begin{center}
                    \textcolor{red}{Quiconque} apprend le français sait que cela prend du temps. \\
                    (Whoever learns French knows that it takes time)
                \end{center}
            \item Personne (nobody), only masculine singular:
                \begin{center}
                    \textcolor{red}{Personne} ne connaît la réponse. \\
                    (Nobody knows the answer)
                \end{center}
            \item Tout (everything), only masculine singular:
                \begin{center}
                    \textcolor{red}{Tout} a l'air en ordre. \\
                    (Everything seems in order)
                \end{center}
            \item Tous (all):
                \begin{center}
                    \begin{tabular}{|c|c|}
                        \hline
                        \textbf{M.P.} & \textbf{F.P.} \\
                        \hline
                        tous & toutes \\
                        \hline
                    \end{tabular}
                \end{center}
                \begin{center}
                    Toutes ont reçu un cadeau pour avoir participé. \\
                    (All got a prize for participating)
                \end{center}
            \item Un, une (one)
        \end{enumerate}

    \section{The Relative Pronoun}
        \begin{enumerate}
            \item \textit{Qui}
                \begin{itemize}
                    \item Replaces the subject which can be a person or a thing.
                    \item Can never be \textit{qu'} in front of a vowel.
                    \item Translate to \textit{who, which, that}.
                    \begin{center}
                        Le café a ouvert l'année dernière. Il a brûlé. \\
                        (The café opened last year. It burned) \\
                        $\to$ Le café \textcolor{red}{qui} a ouvert l'année dernière, a brûlé. \\
                        (The café that opened last year, burned) \\
                        --- \\
                        J'ai vu un film. Ce film était romantique. \\
                        (I saw a movie. This movie was romantic) \\
                        $\to$ J'ai vu un film \textcolor{red}{qui} était romantique. \\
                        (I saw a movie that was romantic)
                    \end{center}
                \end{itemize}
            \item \textit{Qui}
                \begin{itemize}
                    \item Can also be used after the prepositions \textit{à, de, pour, chez, avec}.
                    \item The preposition goes with the verb. In this case, it acts as an indirect object, mostly for people.
                        \begin{center}
                            prêter à \\
                            --- \\
                            Mon ami \textcolor{red}{à qui} j'ai prêté un livre ne me l'a pas rendu. \\
                            (My friend to whom I lent a book didn't give it back to me) \\
                        \end{center}
                \end{itemize}

            \item \textit{Que}
                \begin{itemize}
                    \item Replaces the direct object, becomes \textit{qu'} in front of a vowel or a mute h. 
                    \item Always followed by a pronoun or a noun.
                    \item Translates to \textit{who, whom, which, that}.
                        \begin{center}
                            J'ai vu un film. Le film était nul.
                            $\to$ Le film \textcolor{red}{que} j'ai vu, était nul. \\
                            (The movie that I saw, was bad)
                        \end{center}
                \end{itemize}

            \item \textit{Dont}
                \begin{itemize}
                    \item Replaces an \underline{object} or a \underline{person} when it includes \textit{de}. 
                    \item The preposition \textit{de} goes with the verb.
                        \begin{center}
                            Le jardin \textcolor{red}{dont} je m'occupe est en mauvais état. \\
                            (The garden that I take care of is in bad shape) 
                        \end{center}
                \end{itemize}

            \item \textit{Où}
                \begin{itemize}
                    \item Replaces a \underline{place} or \underline{a place in time}.
                    \item Translates to \textit{where, when}.
                        \begin{center}
                            Le jour \textcolor{red}{où} tu es né était merveilleux. \\
                            (The day you were born was wonderful) \\
                            C'est là \textcolor{red}{où} je travaille. (That's where I work)
                        \end{center}
                \end{itemize}

            \item \textit{Lequel}
                \begin{itemize}
                    \item Used after a preposition other than \textit{de}.
                    \item The four different forms are \textit{lequel, laquelle, lesquels, lesquelles}.
                        \begin{center}
                            C'est la robe \textcolor{red}{avec laquelle} je me suis mariée. \\
                            (It's the dress with which I got married)
                        \end{center}
                \end{itemize}
            
            \item \textit{Auquel = à + lequel}
                \begin{itemize}
                    \item Used for preposition \textit{à}.
                    \item The four different forms are \textit{auquel, à laquelle, auxquels, auxquelles}.
                        \begin{center}
                            Parler à \\
                            --- \\
                            Le client \textcolor{red}{auquel} je parle est impoli / ... \textcolor{red}{à qui je parle} est impoli. \\
                            (The client to whom I am talking is rude)
                        \end{center}
                \end{itemize}

            \item \textit{Duquel = de + lequel}
                \begin{itemize}
                    \item Used to replace a preposition including \textit{de}, but not \textit{de} alone.
                    \item Common ones are \textit{près de, loin de, à côté de, à cause de, à la place de}.
                    \item The four different forms are \textit{duquel, de laquelle, desquels, desquelles}.
                        \begin{center}
                            C'est le bateau \textcolor{red}{au bord duquel} nous avons voyagé. \\
                            (It's the boat on the edge of which we traveled)
                        \end{center}
                \end{itemize}
        \end{enumerate}
    
    \section{Ce}
        A neutral pronoun, can be replaced by \textit{la chose}.
        \begin{enumerate}
            \item \textit{Ce qui}: 
                \begin{center}
                    Tu as vu \textcolor{red}{ce qui} s'est passé? (Did you see what happened?)
                \end{center}
            
            \item \textit{Ce que}:
                \begin{center}
                    Tu comprends \textcolor{red}{ce que} je t'explique (Do you understand what I'm explaining to you?)
                \end{center}
            
            \item \textit{Ce dont}:
                \begin{center}
                    Je sais \textcolor{red}{ce dont} tu as peur. (I know what you are afraid of)
                \end{center}

            \item \textit{Ce $?$ quoi}:
                \begin{itemize}
                    \item à: s'intéresser à
                        \begin{center}
                            Le veux savoir \textcolor{red}{ce à quoi} tu t'interesses. \\
                            (I want to know what you are interested in)
                        \end{center}

                    \item avec: jouner avec
                        \begin{center}
                            Tu ne devineras jamais \textcolor{red}{ce qvec quoi} le chien joue. \\
                            (You will never guess what the dog is playing with)
                        \end{center}

                    \item en: croire en
                        \begin{center}
                            Il parle de \textcolor{red}{ce en quoi} il croit. \\
                            (He talks about what he believes in)
                        \end{center}
                    
                    \item pour: travailler pour
                        \begin{center}
                            \textcolor{red}{Ce pour quoi} il travaille, c'est sa famille. \\
                            (What he works for, is his family)
                        \end{center}
                    
                    \item sur: s'appayer sur
                        \begin{center}
                            \textcolor{red}{Ce sur quoi} je m'appuie, ce sont les lois. \\
                            (What I am relying on, are the laws)
                        \end{center}
                \end{itemize}

            \item \textit{Ce = la chose}
                \begin{center}
                    Tu as vu \textcolor{red}{ce qui} s'est passé $\longleftrightarrow$ Tu as vu \textcolor{red}{la chose qui} s'est passée? \\
                    (Did you see what just happened?)
                \end{center}

            \item \textit{Ce qui}: the subject of the verb
                \begin{center}
                    \textcolor{red}{Ce qui} ne fait peurm c'est le vide. \\
                    (What is scary is the emptiness)
                \end{center}
            
            \item \textit{Ce que}: the direct object, used for verb with no preposition
                \begin{center}
                    C'est \textcolor{red}{ce que} je veux dire. \\
                    (That's what I want to say)
                \end{center}
                
            \item \textit{Ce dont}: when verb is followed by \textit{de}
                \begin{center}
                    \textcolor{red}{Ce dont} elle rêve, c'est de devenir avocate. \\
                    (What she dreams of, is to become a lawyer)
                \end{center}
            
            \item \textit{Ce \dots quoi}: used for all other prepositions except \textit{de}
                \begin{center}
                    \textcolor{red}{Ce à quoi} je pense, c'est mon avenir. \\
                    (What I am thinking about, is my future)
                \end{center}
        \end{enumerate}

\chapter{Les Adverbs}

    Can be one word or a group of words, don't have gender.

    \section{Words ending in -ment}
        \begin{enumerate}
            \item Add \textit{-ment} for vowel ending: \textit{Bizarre} $\to$ \textit{bizarrement} (strangely)
            \item For consonant ending, add \textit{-ment} to the feminine form: \textit{Chaud} $\to$ \textit{chaudement}
                    * Two exceptions: \textit{gentil} $\to$ \textit{gentiment} \& \textit{bref (brève)} $\to$ \textit{brièvement}
            \item Last \textit{e} become \textit{é + ment}: \textit{Aveugle} $\to$ \textit{aveuglément} (blindly)
            \item Adjectives ending in \textit{-ant, -ent}, endings for adverbs are \textit{-amment, -emment}
        \end{enumerate}
        
        \subsection{Irréguliers}
            \begin{itemize}
                \item \textit{Bon} $\to$ \textit{bien}
                \item Mauvais $\to$ mal
                \item Meilleur $\to$ mieux
                \item Petit $\to$ peu
            \end{itemize}
    
    \section{Common Adverbs}
        \begin{enumerate}
            \item For places:
                \begin{center}
                    \begin{tabular}{|c|c|}
                        \hline
                        à droite de & to the right of \\
                        \hline
                        autour de & around \\
                        \hline
                        dehors & outside \\
                        \hline
                        en bas & downstairs \\
                        \hline
                        en haut & upstairs \\
                        \hline
                        \dots & \dots \\
                        \hline
                    \end{tabular}
                \end{center}

            \item Time \& Frequency:
                \begin{center}
                    \begin{tabular}{|c|c|}
                        \hline
                        bientôt & soon \\
                        \hline
                        ensuite & then \\
                        \hline
                        hier & yesterday \\
                        \hline
                        jamais & never \\
                        \hline
                        rarement & rarely \\
                        \hline
                        \dots & \dots \\
                        \hline
                    \end{tabular}
                \end{center}

            \item Quantity:
                \begin{center}
                    \begin{tabular}{|c|c|}
                        \hline
                        assez & enough \\
                        \hline
                        encore & again \\
                        \hline
                        trop & too much \\
                        \hline
                        \dots & \dots \\
                        \hline
                    \end{tabular}
                \end{center}
                
            \item Reason:
                \begin{center}
                    \begin{tabular}{|c|c|}
                        \hline
                        cependaant & however \\
                        \hline
                        donc & so \\
                        \hline
                        même & even \\
                        \hline
                        toutefois & nevertheless \\
                        \hline
                        par conséquent & consequently \\
                        \hline
                        pourtant & yet \\
                        \hline
                        tout de même & all the same \\
                        \hline
                        \dots & \dots \\
                        \hline
                    \end{tabular}
                \end{center}

            \item Manner:
                \begin{center}
                    \begin{tabular}{|c|c|}
                        \hline
                        ainsi & since \\
                        \hline
                        calmement & calmly \\
                        \hline
                        mieux & better \\
                        \hline
                        surtout & mostly \\
                        \hline
                        vite & quickly \\
                        \hline
                        \dots & \dots \\
                        \hline
                    \end{tabular}
                \end{center}

            \item Positive Affirmation:
                \begin{center}
                    \begin{tabular}{|c|c|}
                        \hline
                        peut-être & maybe \\
                        \hline
                        sans doute & without a doubt \\
                        \hline
                        volontiers & willingly \\
                        \hline
                        vraiment & really \\
                        \hline
                        \dots & \dots \\
                        \hline
                    \end{tabular}
                \end{center}

            \item Negative Affirmation:
                \begin{center}
                    \begin{tabular}{|c|c|}
                        \hline
                        ne \dots jamais & never \\
                        \hline
                        pas du tout & not at all \\
                        \hline
                        \dots & \dots \\
                        \hline
                    \end{tabular}
                \end{center}

            \item Question:
                \begin{center}
                    \begin{tabular}{|c|c|}
                        \hline
                        Combien & how much \\
                        \hline
                        \dots & \dots \\
                        \hline
                    \end{tabular}
                \end{center}

            \item Conjunction:
                \begin{center}
                    \begin{tabular}{|c|c|}
                        \hline
                        ainsi & thus \\
                        \hline
                        alors & so \\
                        \hline
                        ensuite & then \\
                        \hline
                        enfin & finally \\
                        \hline
                        par contre & on the other hand \\
                        \hline
                        pourtant & yet \\
                        \hline
                        puis & then \\
                        \hline
                        \dots & \dots \\
                        \hline
                    \end{tabular}
                \end{center}
        \end{enumerate}

    \section{Where to Place Adverb}
        \begin{enumerate}
            \item 90\% of the time: [subject + verb + adv] 
            \item For longer ones, at beginning or end.
                \begin{center}
                    En général, il a de bonnes notes $\leftrightarrow$ Il a de bonnes notes en général. \\
                    (In general, he has good grades)
                \end{center}
            \item Passé Composé, futur prochem passé récent etc: [subject + avoir / être + adv. + past participle]
                \begin{center}
                    \begin{tabular}{|c|c|}
                        \hline
                        Assez & enough \\
                        \hline
                        Bien & well \\
                        \hline
                        Beaucoup & a lot \\
                        \hline
                        bientôt & soon \\
                        \hline
                        Déjà & already \\
                        \hline
                        Encore & again \\
                        \hline
                        Enfin & finally \\
                        \hline
                        Jamais & never \\
                        \hline
                        Mal & badly \\
                        \hline
                        Mieux & better \\
                        \hline
                        Moins & less \\
                        \hline
                        Souvent & often \\
                        \hline
                        Toujours & always \\
                        \hline
                        Trop & too much \\
                        \hline
                        Vite & quickly \\
                        \hline
                    \end{tabular}
                \end{center}
            \item \textit{Trop} and also be followed by other adverb:
                \begin{center}
                    \begin{tabular}{|c|c|}
                        \hline
                        Trop souvent & too often \\
                        \hline
                        Trop vite & too quickly \\
                        \hline
                        Trop mal & too badly \\
                        \hline
                        \dots & \dots \\
                        \hline
                    \end{tabular}
                \end{center}
            \item Adverbs can be anywhere: \textit{hier}
        \end{enumerate}

    \section{Bon \& Bien}
        \begin{enumerate}
            \item \textit{Bon} can be used as adjective, adverb, noun.
                \begin{itemize}
                    \item Adjective: \textit{bon, bonne, bons, bonnes} (good, nice, right)
                    \item Adverb: \textit{bon}
                    \item Noun: \textit{bon}, and \textit{un bon \textcolor{red}{d'achat}} (a voucher / a coupon), the text highlighted in red can be omitted.
                    \item In expression (adj):
                        \begin{center}
                            \begin{tabular}{|c|c|}
                                \hline
                                pour de bon & for good \\
                                \hline
                                bon marché & cheap \\
                                \hline
                                \dots & \dots \\
                                \hline
                            \end{tabular}
                        \end{center}
                    \item In greeting \& wishes (adj): \textit{Bonne année} (Happy New Year)
                \end{itemize}
            
            \item \textit{Bien} can be adjective, adverb, noun, always stay as \textit{bien}.
                \begin{itemize}
                    \item Good: \textit{C'est bien. Très bien! Ça a l'air bien} (it seems good)
                    \item Verb + bien: \textit{Tu vas bien? Je sais bien. Le wifi ne fonctionne pas bien}
                    \item Auxiliary + bien + verb: \textit{Nous avons \textcolor{red}{bien} étudié} (We studied well)
                    \item Noun: \textit{un bien \textcolor{red}{immobiliér}}, text in red can be omitted.
                    \item In expressions:
                        \begin{center}
                            Elle est \textcolor{red}{bien} dans sa peau. (She is comfortable in her skin) \\
                            Être \textcolor{red}{bien} arrivé(e) (to arrive safely) \\
                            \vspace{0.2cm}
                            \begin{tabular}{|c|c|}
                                \hline
                                Aller bien & to be well \\
                                \hline
                                Bien entendu & of course \\
                                \hline
                                Bien sûr que [non $\smile$ qui $\smile$ $\emptyset$] & of course not / yes, of course / of course \\
                                \hline
                                on verra bien & we will see \\
                                \hline
                                bien prendre & to take something well \\
                                \hline
                                bien souvent & very often \\
                                \hline
                                faire bien de & to do the right / good thing \\
                                \hline
                                être bien se sentir bien dans sa peau & to feel comfortable in ones skin \\
                                \hline
                            \end{tabular}
                        \end{center}
                \end{itemize}
            
            \item \textit{Bon} vs. \textit{Bien}
                \begin{itemize}
                    \item Use \textit{C'est bon} for physical sensation, taste. Note that gender rule doesn't apply after \textit{ce}.
                    \item Use \textit{C'est bien} for judgement, opinion.
                \end{itemize}
        \end{enumerate}

    \section{Encore \& Toujours}
        \begin{itemize}
            \item Translate to again \& always, both can be still or yet
            \item \textit{Toujours} can also mean anyhow / anyway
            \item \textit{Encore} can mean more / another / even
        \end{itemize}

    \section{Déjà \& Jamais}
        \begin{itemize}
            \item \textit{Déjà} can be yet in a positive question.
                \begin{center}
                    Tu as \textcolor{red}{déjà} appris ça? (Have you learned this yet?)
                \end{center}
            \item \textit{Déjà} can be even in question have you ever \dots
                \begin{center}
                    Tu as \textcolor{red}{déjà} été à Paris? (Have you ever been to Paris?)
                \end{center}
            \item \textit{Jamais} can be never in question have you never \dots
                \begin{center}
                    Tu n'as jamais vu ce film? (Have you never seen this movie?)
                \end{center}
        \end{itemize}

\chapter{Conjunction}

    \section{Coordinating Conjunction}
        \begin{center}
            \begin{tabular}{|c|c|}
                \hline
                mais & but \\
                \hline
                ou & or \\
                \hline
                et & and \\
                \hline
                danc & so \\
                \hline
                or & yet \\
                \hline
                ni & nor \\
                \hline
                car & because \\
                \hline
            \end{tabular}
        \end{center}

    \section{Subordinating Conjunctions}
        \begin{center}
            \begin{tabular}{|c|c|}
                \hline
                Comme & as / since, to compare \\
                \hline
                Puisque & as / since \\
                \hline
                Quand & when \\
                \hline
                Lorsque & when \\
                \hline
                que & that \\
                \hline
                quoique & although \\
                \hline
                si & if \\
                \hline
            \end{tabular}
        \end{center}

    \section{Other Conjunctions}
        \begin{center}
            \begin{tabular}{|c|c|}
                \hline
                à moin que & unless that \\
                \hline
                de peur que & for fear that \\
                \hline
                depuis que & since \\
                \hline
                même si & even if \\
                \hline
                \dots & \dots \\
                \hline
            \end{tabular}
        \end{center}
\end{document}